% Generated mechanically from frozen input p12/ko/moduli.tex.
% Do not edit this generated file; edit the bound locale source instead.

% OK, start here.
%

\setcounter{chapter}{107}
\chapter{모듈라이 스택}


\phantomsection
\label{moduli-section-phantom}





\section{서론}
\label{moduli-section-introduction}

\noindent
이 장에서는
$\mathit{Hom}$, $\mathit{Isom}$, $\Cohstack_{X/B}$,
$\Quotfunctor_{\mathcal{F}/X/B}$, $\Hilbfunctor_{X/B}$,
$\Picardstack_{X/B}$, $\Picardfunctor_{X/B}$, $\mathit{Mor}_B(Z, X)$,
$\Spacesstack'_{fp, flat, proper}$, $\Polarizedstack$, 그리고
$\Complexesstack_{X/B}$와 같은 모듈라이 공간과 모듈라이 스택의
기본 성질을 확인한다.
적절한 가정 아래에서 이들이 대수공간 또는 대수적 스택임은 이미 보였다.
『몫』의 절
\ref{quot-section-hom},
\ref{quot-section-isom},
\ref{quot-section-stack-coherent-sheaves},
\ref{quot-section-not-flat},
\ref{quot-section-quot},
\ref{quot-section-hilb},
\ref{quot-section-picard-stack},
\ref{quot-section-picard-functor},
\ref{quot-section-relative-morphisms},
\ref{quot-section-stack-of-spaces},
\ref{quot-section-polarized}, 그리고
\ref{quot-section-moduli-complexes}를 보라.
곡선의 스택은 $\textit{Curves}$로 표기하며 『몫』의 절
\ref{quot-section-curves}에서 도입했다. 이에 대해서는 곡선의 모듈라이를
다루는 장, 즉 『곡선의 모듈라이』의 절
\ref{moduli-curves-section-stack-curves}에서 논한다.

\medskip\noindent
어떤 의미에서 이 장은 Grothendieck의 강의
\cite{Gr-I},
\cite{Gr-II},
\cite{Gr-III},
\cite{Gr-IV},
\cite{Gr-V}, 그리고
\cite{Gr-VI}의 발자취를 따른다.







\section{규약과 용어의 관용적 사용}
\label{moduli-section-conventions}

\noindent
『스택의 성질』의 절 \ref{stacks-properties-section-conventions}에서
도입한 규약과 용어의 관용적 사용을 계속 따른다.
달리 언급하지 않는 한 바탕 스킴은 $\Spec(\mathbf{Z})$이다.






\section{Hom과 Isom의 성질}
\label{moduli-section-hom-isom}

\noindent
$f : X \to B$를 유한 표시인 대수공간의 사상이라 하자.
$\mathcal{F}$와 $\mathcal{G}$가 준연접 $\mathcal{O}_X$-가군이라고 하자.
$\mathcal{G}$가 유한 표시이고 $B$ 위에서 평탄하며 그 지지가
$B$ 위에서 고유이면, 다음으로 정의되는 함자
$\mathit{Hom}(\mathcal{F}, \mathcal{G})$
$$
T/B \longmapsto \Hom_{\mathcal{O}_{X_T}}(\mathcal{F}_T, \mathcal{G}_T)
$$
는 $B$ 위의 아핀 대수공간이다. $\mathcal{F}$가 유한 표시이면
$\mathit{Hom}(\mathcal{F}, \mathcal{G}) \to B$는 유한 표시이다.
『몫』의 명제 \ref{quot-proposition-hom}을 보라.

\medskip\noindent
$\mathcal{F}$와 $\mathcal{G}$가 모두 유한 표시이고 $B$ 위에서 평탄하며
그 지지가 $B$ 위에서 고유이면, 부분함자
$$
\mathit{Isom}(\mathcal{F}, \mathcal{G}) \subset
\mathit{Hom}(\mathcal{F}, \mathcal{G})
$$
는 $B$ 위에서 유한 표시인 아핀 대수공간이다.
『몫』의 명제 \ref{quot-proposition-isom}을 보라.




\section{연접층의 스택의 성질}
\label{moduli-section-stack-coherent-sheaves}

\noindent
$f : X \to B$를 분리이고 유한 표시인 대수공간의 사상이라 하자.
그러면 고유 지지를 갖는 연접 가군의 평탄한 족을 매개화하는 스택
$\Cohstack_{X/B}$는 대수적이다. 『몫』의 정리
\ref{quot-theorem-coherent-algebraic-general}를 보라.

\begin{lemma}
\label{moduli-lemma-coherent-diagonal-affine-fp}
$B$ 위에서 $\Cohstack_{X/B}$의 대각사상은 아핀이며 유한 표시이다.
\end{lemma}

\begin{proof}
대각사상이 대수공간으로 표현가능함은 『몫』의 보조정리
\ref{quot-lemma-coherent-diagonal}에서 보였다.
그 증명을 살펴보면, $T$ 위에서 평탄하고 지지가 $T$ 위에서 고유인
유한 표시 $\mathcal{O}_{X_T}$-가군의 쌍
$\mathcal{F}$, $\mathcal{G}$에 대해
$\mathit{Isom}(\mathcal{F}, \mathcal{G}) \to T$가 아핀이며
유한 표시임을 보이면 된다. 이는 절 \ref{moduli-section-hom-isom}에서 논했다.
\end{proof}

\begin{lemma}
\label{moduli-lemma-coherent-qs-lfp}
사상 $\Cohstack_{X/B} \to B$는 준분리이고 국소적으로 유한 표시이다.
\end{lemma}

\begin{proof}
$\Cohstack_{X/B} \to B$가 준분리임을 확인하려면 그 대각사상이
준콤팩트이고 준분리임을 보여야 한다. 이는 보조정리
\ref{moduli-lemma-coherent-diagonal-affine-fp}에서 즉시 따른다.
$\Cohstack_{X/B} \to B$가 국소적으로 유한 표시임을 보이려면
$\Cohstack_{X/B} \to B$가 극한을 보존함을 보이면 된다.
『스택의 극한』의 명제
\ref{stacks-limits-proposition-characterize-locally-finite-presentation}를 보라.
이는 『몫』의 보조정리 \ref{quot-lemma-coherent-limits}에서 따른다
(작은 세부사항은 생략한다).
\end{proof}

\begin{lemma}
\label{moduli-lemma-coherent-existence-part}
$X \to B$가 유한 표시일 뿐 아니라 고유라고 하자.
그러면 $\Cohstack_{X/B} \to B$는 값매김 판정법의 존재 부분을 만족한다
(『스택의 사상』의 정의 \ref{stacks-morphisms-definition-existence}).
\end{lemma}

\begin{proof}
밑변환을 취하면 곧바로 다음 문제로 귀착된다. 분수체가 $K$인
값매김환 $R$, $R$ 위에서 고유인 대수공간 $X$, 그리고 연접
$\mathcal{O}_{X_K}$-가군 $\mathcal{F}_K$가 주어졌을 때, 일반 섬유가
$\mathcal{F}_K$인, $R$ 위에서 평탄한 유한 표시
$\mathcal{O}_X$-가군 $\mathcal{F}$가 존재함을 보여라.
『대수공간 위의 평탄성』의 정리
\ref{spaces-flat-theorem-finite-type-flat}에 의해 $R$ 위에서 평탄한
유한형 준연접 $\mathcal{O}_X$-가군 $\mathcal{F}$는 모두 유한 표시임에
유의하자. 일반 섬유의 매장을 $j : X_K \to X$로 표기하자.
아핀 사상 $\Spec(K) \to \Spec(R)$의 밑변환이므로 $j$는 아핀이다.
따라서 $j_*\mathcal{F}_K$는 준연접이다. 이를 유한형 준연접
$\mathcal{O}_X$-부분가군들의 여과 쌍대극한
$$
j_*\mathcal{F}_K = \colim \mathcal{F}_i
$$
으로 쓰자. 『대수공간의 극한』의 보조정리
\ref{spaces-limits-lemma-directed-colimit-finite-type}를 보라.
$j_*\mathcal{F}_K$는 $X$ 위의 $K$-벡터공간 층이므로
$\Spec(R)$ 위에서 평탄하다. 값매김환 위에서 평탄하다는 것은
꼬임이 없다는 것과 같고
(『대수학 더 보기』의 보조정리
\ref{more-algebra-lemma-valuation-ring-torsion-free-flat}),
부분가군은 꼬임 없음 성질을 물려받으므로 각 $\mathcal{F}_i$는
$R$ 위에서 평탄하다.
마지막으로 어떤 $i$에 대해 사상
$j^*\mathcal{F}_i \to \mathcal{F}_K$가 동형임을 보여야 한다.
$j^*j_*\mathcal{F}_K = \mathcal{F}_K$이고(작은 세부사항은 생략한다)
$j^*$는 완전하므로, 모든 $i$에 대해
$j^*\mathcal{F}_i \to \mathcal{F}_K$가 단사임을 알 수 있다.
$j^*$는 쌍대극한과 가환하므로
$\mathcal{F}_K = j^*j_*\mathcal{F}_K = \colim j^*\mathcal{F}_i$이다.
$\mathcal{F}_K$는 연접, 즉 유한 표시이므로 어떤 $i$가 존재하여
$j^*\mathcal{F}_i$가 $X$의 어떤 아핀 에탈 덮개 위에서 모든
(유한 개의) 생성원을 포함한다.
따라서 충분히 큰 $i$에 대해
$j^*\mathcal{F}_i \to \mathcal{F}_K$가 전사임을 얻는다.
\end{proof}

\begin{lemma}
\label{moduli-lemma-coherent-functorial}
$B$를 대수공간이라 하자. $\pi : X \to Y$를 $B$ 위에서 분리이고
유한 표시인 대수공간들 사이의 준유한 사상이라 하자. 그러면
$\pi_*$는 사상 $\Cohstack_{X/B} \to \Cohstack_{Y/B}$를 유도한다.
\end{lemma}

\begin{proof}
$(T \to B, \mathcal{F})$를 $\Cohstack_{X/B}$의 대상이라 하자.
다음을 주장한다.
\begin{enumerate}
\item[(a)] $(T \to B, \pi_{T, *}\mathcal{F})$는
$\Cohstack_{Y/B}$의 대상이고,
\item[(b)] $T' \to T$에 대해
$\pi_{T', *}(X_{T'} \to X_T)^*\mathcal{F} =
(Y_{T'} \to Y_T)^*\pi_{T, *}\mathcal{F}$이다.
\end{enumerate}
(b)는 이 구성이 원하는 함자
$\Cohstack_{X/B} \to \Cohstack_{Y/B}$를 정의함을 보장한다.

\medskip\noindent
$\mathcal{F}$의 영 번째 맞춤 아이디얼로 잘라낸 닫힌 부분공간을
$i : Z \to X_T$라 하자
(『대수공간 위의 인자』의 절
\ref{spaces-divisors-section-fitting-ideals}).
그러면 가정에 의해 $Z \to B$는 고유이다(『대수공간의 유도 범주』의 절
\ref{spaces-perfect-section-proper-over-base}을 보라).
한편 $i$는 유한 표시이다
(『대수공간 위의 인자』의 보조정리
\ref{spaces-divisors-lemma-fitting-ideal-of-finitely-presented}와
『대수공간의 사상』의 보조정리
\ref{spaces-morphisms-lemma-closed-immersion-finite-presentation}).
$i_*\mathcal{G} = \mathcal{F}$를 만족하는 유한형 준연접
$\mathcal{O}_Z$-가군 $\mathcal{G}$가 존재한다
(『대수공간 위의 인자』의 보조정리
\ref{spaces-divisors-lemma-on-subscheme-cut-out-by-Fit-0}).
실제로 『대수공간 위의 내림』의 보조정리
\ref{spaces-descent-lemma-finite-finitely-presented-module}에 의해
$\mathcal{G}$는 $\mathcal{O}_Z$-가군으로서 유한 표시이다.
예를 들어 $\mathcal{G}$와 $\mathcal{F}$의 줄기가 일치하므로
(『대수공간의 사상』의 보조정리
\ref{spaces-morphisms-lemma-stalk-push-closed}),
$\mathcal{G}$는 $B$ 위에서 평탄하다.
$\pi_T \circ i : Z \to Y_T$는 준유한 사상들의 합성이므로 준유한이고
$\pi_{T, *}\mathcal{F} = (\pi_T \circ i)_*\mathcal{G})$임에 유의하자.
$i$는 아핀이므로 $i_*$의 형성은 밑변환과 가환한다
(『대수공간의 코호몰로지』의 보조정리
\ref{spaces-cohomology-lemma-affine-base-change}).
따라서 $B$를 $T$로, $X$를 $Z$로, $\mathcal{F}$를 $\mathcal{G}$로,
$Y$를 $Y_T$로 바꾸어 다음 문단에서 다루는 경우로 귀착시킬 수 있다.

\medskip\noindent
$X \to B$가 고유라고 하자. 그러면 『대수공간의 사상』의 보조정리
\ref{spaces-morphisms-lemma-universally-closed-permanence}에 의해
$\pi$는 고유이고, 따라서 『대수공간의 사상 더 보기』의 보조정리
\ref{spaces-more-morphisms-lemma-characterize-finite}에 의해 유한이다.
유한 사상은 아핀이므로 『대수공간의 코호몰로지』의 보조정리
\ref{spaces-cohomology-lemma-affine-base-change}에 의해 (b)가 성립한다.
한편 『대수공간의 사상』의 보조정리
\ref{spaces-morphisms-lemma-finite-presentation-permanence}에 의해
$\pi$는 유한 표시이다. 따라서 『대수공간 위의 내림』의 보조정리
\ref{spaces-descent-lemma-finite-finitely-presented-module}에 의해
$\pi_{T, *}\mathcal{F}$는 유한 표시이다. 마지막으로 예를 들어
『대수공간의 코호몰로지』의 보조정리
\ref{spaces-cohomology-lemma-stalk-push-finite}를 사용하여 줄기를 보면
$\pi_{T, *}\mathcal{F} $는 $B$ 위에서 평탄하다.
\end{proof}

\begin{lemma}
\label{moduli-lemma-coherent-open}
$B$를 대수공간이라 하자. $\pi : X \to Y$를 $B$ 위에서 분리이고
유한 표시인 대수공간들의 열린 몰입이라 하자. 그러면 보조정리
\ref{moduli-lemma-coherent-functorial}의 사상
$\Cohstack_{X/B} \to \Cohstack_{Y/B}$는 열린 몰입이다.
\end{lemma}

\begin{proof}
생략한다. 도움말: $\mathcal{F}$가 $T$ 위의 $\Cohstack_{Y/B}$의 대상이고
$t \in T$에 대해 $\text{Supp}(\mathcal{F}_t) \subset |X_t|$이면,
$t$의 어떤 근방에 있는 모든 $t' \in T$에 대해서도 같은 명제가 성립한다.
\end{proof}

\begin{lemma}
\label{moduli-lemma-coherent-closed}
$B$를 대수공간이라 하자. $\pi : X \to Y$를 $B$ 위에서 분리이고
유한 표시인 대수공간들의 닫힌 몰입이라 하자. 그러면 보조정리
\ref{moduli-lemma-coherent-functorial}의 사상
$\Cohstack_{X/B} \to \Cohstack_{Y/B}$는 닫힌 몰입이다.
\end{lemma}

\begin{proof}
$X$를 $Y$의 닫힌 부분공간으로 잘라내는 아이디얼 층을
$\mathcal{I} \subset \mathcal{O}_Y$라 하자. $\pi_*$가 준연접
$\mathcal{O}_X$-가군의 범주와 $\mathcal{I}$에 의해 소멸되는 준연접
$\mathcal{O}_Y$-가군의 범주 사이의 동치를 유도함을 상기하자.
『대수공간의 사상』의 보조정리
\ref{spaces-morphisms-lemma-i-star-equivalence}를 보라.
$T \to B$로 밑변환한 뒤에도 필요한 수정을 가하면 같은 명제가 성립한다.
이때 $\mathcal{I}$는 아이디얼 층
$\mathcal{I}_T = \Im((Y_T \to Y)^*\mathcal{I} \to \mathcal{O}_{Y_T})$로
대체된다. 보조정리 \ref{moduli-lemma-coherent-functorial}의 증명을 분석하면
$\Cohstack_{X/B} \to \Cohstack_{Y/B}$의 본질적 상은 정확히
$\mathcal{F}$가 $\mathcal{I}_T$에 의해 소멸되는 대상
$\xi = (T \to B, \mathcal{F})$임을 알 수 있다.
다시 말해, $\xi$가 본질적 상에 속할 필요충분조건은 곱셈 사상
$$
\mathcal{F} \otimes_{\mathcal{O}_{Y_T}} (Y_T \to Y)^*\mathcal{I}
\longrightarrow
\mathcal{F}
$$
이 영이고, 임의의 추가 밑변환 $T' \to T$ 뒤에도 마찬가지인 것이다.
다음에 유의하자.
$$
(Y_{T'} \to Y_T)^*(
\mathcal{F} \otimes_{\mathcal{O}_{Y_T}} (Y_T \to Y)^*\mathcal{I}) =
(Y_{T'} \to Y_T)^*\mathcal{F} \otimes_{\mathcal{O}_{Y_{T'}}}
(Y_{T'} \to Y)^*\mathcal{I})
$$
따라서 $T'$ 위에서 곱셈 사상이 영이 되는 조건은
『대수공간 위의 평탄성』의 보조정리
\ref{spaces-flat-lemma-F-zero-closed-proper}에 의해 $T$의 닫힌 부분공간으로
표현가능하다.
\end{proof}

\begin{situation}[수치 불변량]
\label{moduli-situation-numerical}
$f : X \to B$를 이 절의 도입부에서와 같이 두자. $I$를 집합이라 하고,
각 $i \in I$에 대해 $E_i \in D(\mathcal{O}_X)$가 완전하다고 하자.
$\Cohstack_{X/B}$의 대상 $(T \to B, \mathcal{F})$가 주어졌을 때,
$X_T$로의 $E_i$의 유도 당김을 $E_{i, T}$로 표기하자.
$D(\mathcal{O}_T)$의 대상
$$
K_i = Rf_{T, *}(E_{i, T} \otimes_{\mathcal{O}_{X_T}}^\mathbf{L} \mathcal{F})
$$
는 완전하고 그 형성은 밑변환과 가환한다. 『대수공간의 유도 범주』의
보조정리 \ref{spaces-perfect-lemma-base-change-tensor-perfect}를 보라.
따라서 함수
$$
\chi_i : |T| \longrightarrow \mathbf{Z},\quad
\chi_i(t) =
\chi(X_t, E_{i, t} \otimes_{\mathcal{O}_{X_t}}^\mathbf{L} \mathcal{F}_t) =
\chi(K_i \otimes_{\mathcal{O}_T}^\mathbf{L} \kappa(t))
$$
는 『대수공간의 유도 범주』의 보조정리
\ref{spaces-perfect-lemma-chi-locally-constant}에 의해 국소 상수이다.
$P : I \to \mathbf{Z}$를 사상이라 하자. 수치 불변량이 $P$와 일치하는,
고유 지지를 갖는 연접층의 평탄한 족으로 이루어진 부분스택
$$
\Cohstack^P_{X/B} \subset \Cohstack_{X/B}
$$
을 생각하자. 더 정확히 말해 $\Cohstack_{X/B}$의 대상
$(T \to B, \mathcal{F})$가 $\Cohstack^P_{X/B}$에 속할 필요충분조건은
모든 $i \in I$와 $t \in T$에 대해 $\chi_i(t) = P(i)$인 것이다.
\end{situation}

\begin{lemma}
\label{moduli-lemma-open-P}
상황 \ref{moduli-situation-numerical}에서 스택 $\Cohstack^P_{X/B}$는 대수적이고
$$
\Cohstack^P_{X/B} \longrightarrow \Cohstack_{X/B}
$$
는 평탄한 닫힌 몰입이다. $I$가 유한이거나 $B$가 국소 뇌터이면,
$\Cohstack^P_{X/B}$는 $\Cohstack_{X/B}$의 열린 동시에 닫힌 부분스택이다.
\end{lemma}

\begin{proof}
$I$가 유한이면 함수 $t \mapsto \chi_i(t)$들이 국소 상수이므로 이는
즉시 명백하다. $I$가 무한이면
$$
I = \bigcup\nolimits_{I' \subset I\text{ finite}} I'
$$
로 쓰고 $P' = P|_{I'}$로 표기한다. 그러면
$$
\Cohstack^P_{X/B} = \bigcap\nolimits_{I' \subset I\text{ finite}}
\Cohstack^{P'}_{X/B}
$$
이다. 따라서 $\Cohstack^P_{X/B}$는 항상 대수적 스택이고 사상
$\Cohstack^P_{X/B} \subset \Cohstack_{X/B}$는 항상 평탄한 닫힌
몰입이지만, 더 이상 열린 부분스택이 아닐 수도 있다(예를 구성하는 일은
독자에게 맡긴다). 하지만 $B$가 국소 뇌터이면 보조정리
\ref{moduli-lemma-coherent-qs-lfp}와 『스택의 사상』의 보조정리
\ref{stacks-morphisms-lemma-locally-finite-type-locally-noetherian}에 의해
$\Cohstack_{X/B}$도 국소 뇌터이다. 따라서 $U$가 국소 뇌터 스킴이고
$U \to \Cohstack_{X/B}$가 매끄러운 전사 사상이면, 열린 동시에 닫힌
부분스택 $\Cohstack^{P'}_{X/B}$들의 역상은 $U$에서 열린 교집합을 갖는다
(국소 뇌터 위상공간의 연결 성분은 열려 있기 때문이다).
이로써 이 경우의 결과를 얻는다.
\end{proof}

\begin{lemma}
\label{moduli-lemma-finite-list-perfect-objects}
$f : X \to B$를 이 절의 도입부에서와 같이 두자.
$E_1, \ldots, E_r \in D(\mathcal{O}_X)$가 완전하다고 하자.
$I = \mathbf{Z}^{\oplus r}$라 하고 다음 사상을 생각하자.
$$
I \longrightarrow D(\mathcal{O}_X),\quad
(n_1, \ldots, n_r) \longmapsto
E_1^{\otimes n_1}
\otimes \ldots \otimes
E_r^{\otimes n_r}
$$
$P : I \to \mathbf{Z}$를 사상이라 하자. 상황
\ref{moduli-situation-numerical}에서 정의한
$\Cohstack^P_{X/B} \subset \Cohstack_{X/B}$는 열린 동시에 닫힌
부분스택이다.
\end{lemma}

\begin{proof}
$B$ 위에서 에탈 국소적으로 작업해도 되므로 $B$가 아핀이라고 가정할 수
있다. 이 경우 절대 뇌터 축소를 수행할 수 있다. 독자는 이 증명을
건너뛰어도 좋다. 구체적으로 $B = \Spec(\Lambda)$라 하자.
$\Lambda = \colim \Lambda_i$를 각 $\Lambda_i$가 $\mathbf{Z}$ 위에서
유한형인 여과 쌍대극한으로 쓰자. 어떤 $i$에 대해 분리이고 유한 표시인
대수공간의 사상 $X_i \to \Spec(\Lambda_i)$를 찾아, 이를 $\Lambda$로
밑변환하면 $X$가 되게 할 수 있다. 『대수공간의 극한』의 보조정리
\ref{spaces-limits-lemma-descend-finite-presentation}과
\ref{spaces-limits-lemma-descend-separated-morphism}을 보라.
그런 다음 $i$를 키우면, $D(\mathcal{O}_{X_i})$의 완전 대상
$E_{1, i}, \ldots, E_{r, i}$가 존재하여 $X$로 유도 당김한 것이 각각
$E_1, \ldots, E_r$과 동형이라고 가정할 수 있다.
『대수공간의 유도 범주』의 보조정리
\ref{spaces-perfect-lemma-perfect-on-limit}를 보라.
명백히 다음 정사각형은 카르테시안이다.
$$
\xymatrix{
\Cohstack^P_{X/B} \ar[r] \ar[d] &
\Cohstack_{X/B} \ar[d] \\
\Cohstack^P_{X_i/\Spec(\Lambda_i)} \ar[r] &
\Cohstack_{X_i/\Spec(\Lambda_i)}
}
$$
따라서 보조정리 \ref{moduli-lemma-open-P}를 적용하면 증명이 끝난다.
\end{proof}

\begin{example}[힐베르트 다항식이 고정된 연접층]
\label{moduli-example-hilbert-polynomial}
$f : X \to B$를 이 절의 도입부에서와 같이 두자.
$\mathcal{L}$을 가역 $\mathcal{O}_X$-가군이라 하고,
$P : \mathbf{Z} \to \mathbf{Z}$를 수치 다항식이라 하자.
그러면 수치 불변량이 $P$와 일치하는, 고유 지지를 갖는 연접층의
평탄한 족으로 이루어진 열린 동시에 닫힌 대수적 부분스택
$$
\Cohstack^P_{X/B} =
\Cohstack^{P, \mathcal{L}}_{X/B}
\subset \Cohstack_{X/B}
$$
을 생각할 수 있다. $\Cohstack_{X/B}$의 대상
$(T \to B, \mathcal{F})$가 $\Cohstack^P_{X/B}$에 속할 필요충분조건은
모든 $n \in \mathbf{Z}$와 $t \in T$에 대해
$$
P(n) =
\chi(X_t, \mathcal{F}_t \otimes_{\mathcal{O}_{X_t}} \mathcal{L}_t^{\otimes n})
$$
인 것이다. 물론 이는 $I = \mathbf{Z} \to D(\mathcal{O}_X)$가
$n \mapsto \mathcal{L}^{\otimes n}$으로 주어진 상황
\ref{moduli-situation-numerical}의 특수한 경우이다. 보조정리
\ref{moduli-lemma-finite-list-perfect-objects}에서 이것이 열린 동시에 닫힌
부분스택임이 따른다. 함수
$n \mapsto
\chi(X_t, \mathcal{F}_t \otimes_{\mathcal{O}_{X_t}} \mathcal{L}_t^{\otimes n})$
는 항상 수치 다항식이므로(『체 위의 대수공간』의 보조정리
\ref{spaces-over-fields-lemma-numerical-polynomial-from-euler}),
$$
\Cohstack_{X/B} = \coprod\nolimits_{P\text{ numerical polynomial}}
\Cohstack^P_{X/B}
$$
가 서로소 합 분해임을 얻는다.
\end{example}







\section{Quot의 성질}
\label{moduli-section-quot}

\noindent
$f : X \to B$를 분리이고 유한 표시인 대수공간의 사상이라 하자.
$\mathcal{F}$를 준연접 $\mathcal{O}_X$-가군이라 하자. 그러면
$\Quotfunctor_{\mathcal{F}/X/B}$는 대수공간이다.
$\mathcal{F}$가 유한 표시이면
$\Quotfunctor_{\mathcal{F}/X/B} \to B$는 국소적으로 유한 표시이다.
『몫』의 명제 \ref{quot-proposition-quot}를 보라.

\begin{lemma}
\label{moduli-lemma-quot-diagonal-closed}
$\Quotfunctor_{\mathcal{F}/X/B} \to B$의 대각사상은 닫힌 몰입이다.
$\mathcal{F}$가 유한형이면 이 대각사상은 유한 표시인 닫힌 몰입이다.
\end{lemma}

\begin{proof}
스킴 $T/B$와 $B$ 위의 $\Quotfunctor_{\mathcal{F}/X/B}$의
$T$-값 점들에 대응하는 두 몫
$\mathcal{F}_T \to \mathcal{Q}_i$, $i = 1, 2$가 주어졌다고 하자.
첫 번째 사상의 핵을 $\mathcal{K}_1$로 표기하고, 합성 사상을
$u : \mathcal{K}_1 \to \mathcal{Q}_2$라 두자.
『대수공간 위의 평탄성』의 보조정리
\ref{spaces-flat-lemma-F-zero-closed-proper}에 의해 $T$의 닫힌 부분공간이
존재하여, $T' \to T$가 이를 통해 분해될 필요충분조건은 당김
$u_{T'}$가 영인 것이다. 이는 대각사상이 닫힌 몰입임을 증명한다.
더 나아가 $\mathcal{F}$가 유한형이면 $\mathcal{K}_1$도 유한형이고
(『사이트 위의 가군』의 보조정리
\ref{sites-modules-lemma-kernel-surjection-finite-onto-finite-presentation}),
같은 보조정리에 의해 대각사상이 유한 표시임을 알 수 있다.
\end{proof}

\begin{lemma}
\label{moduli-lemma-quot-s-lfp}
사상 $\Quotfunctor_{\mathcal{F}/X/B} \to B$는 분리이다.
$\mathcal{F}$가 유한 표시이면 이 사상은 국소적으로도 유한 표시이다.
\end{lemma}

\begin{proof}
$\Quotfunctor_{\mathcal{F}/X/B} \to B$가 분리임을 확인하려면
그 대각사상이 닫힌 몰입임을 보여야 한다. 이는 보조정리
\ref{moduli-lemma-quot-diagonal-closed}에 의해 참이다. 두 번째 명제는
『몫』의 명제 \ref{quot-proposition-quot}의 일부이다.
\end{proof}

\begin{lemma}
\label{moduli-lemma-quot-existence-part}
$X \to B$가 유한 표시일 뿐 아니라 고유이고 $\mathcal{F}$가 유한형
준연접 가군이라고 하자. 그러면
$\Quotfunctor_{\mathcal{F}/X/B} \to B$는 값매김 판정법의 존재 부분을
만족한다(『대수공간의 사상』의 정의
\ref{spaces-morphisms-definition-valuative-criterion}).
\end{lemma}

\begin{proof}
밑변환을 취하면 곧바로 다음 문제로 귀착된다. 분수체가 $K$인
값매김환 $R$, $R$ 위에서 고유인 대수공간 $X$, 유한형 준연접
$\mathcal{O}_X$-가군 $\mathcal{F}$, 그리고 연접 몫
$\mathcal{F}_K \to \mathcal{Q}_K$가 주어졌을 때,
$R$ 위에서 평탄하고 일반 섬유가 $\mathcal{Q}_K$인 유한 표시
$\mathcal{O}_X$-가군 $\mathcal{Q}$로 가는 몫
$\mathcal{F} \to \mathcal{Q}$가 존재함을 보여라.
『대수공간 위의 평탄성』의 정리
\ref{spaces-flat-theorem-finite-type-flat}에 의해 $R$ 위에서 평탄한
유한형 준연접 $\mathcal{O}_X$-가군 $\mathcal{F}$는 모두 유한 표시임에
유의하자.
먼저 $\mathcal{Q}$의 존재를 아핀 국소적으로 해결한다.

\medskip\noindent
아핀 국소적으로 다음 문제에 이른다.
$R \to A$를 유한 표시인 환 준동형, $M$을 유한 $A$-가군,
$\varphi : M_K \to N_K$를 $A_K$-몫가군이라 하자. 그러면
$$
L = \{x \in M \mid \varphi(x \otimes 1) = 0 \}
$$
을 생각할 수 있다. $M \to M/L$은 $A$-가군의 몫이고, 이 몫은
$R$-가군으로서 꼬임이 없다. 따라서 $R$-가군으로서 평탄하다
(『대수학 더 보기』의 보조정리
\ref{more-algebra-lemma-valuation-ring-torsion-free-flat}).
$M$이 유한 $A$-가군이므로 $L$도 그러하며, 위의 참고문헌에 의해
$L$은 유한 표시 $A$-가군이라고 결론 내린다. 명백히 $M/L$은
$(M/L)_K = N_K$를 만족하는 이러한 몫으로서 유일하다.

\medskip\noindent
앞 문단의 구성에서의 유일성은 이 몫들이 이어붙여져 원하는
$\mathcal{Q}$를 줌을 보장한다. 조금 더 자세히 설명하자.
$U$가 아핀 스킴인 전사 에탈 사상 $U \to X$를 택하자.
위의 구성을 써서 준연접이고 $R$ 위에서 평탄하며 일반 섬유에서
$\mathcal{Q}_K|U$를 되찾는 몫
$\mathcal{F}|_U \to \mathcal{Q}_U$를 구성한다.
$X$가 분리이므로 $U \times_X U$ 역시 $X$ 위에서 에탈인 아핀
스킴이다. 그러면 구성의 유일성에 의해 몫
$\mathcal{F}|_{U \times_X U} \to \text{pr}_1^*\mathcal{Q}_U$와
$\mathcal{F}|_{U \times_X U} \to \text{pr}_2^*\mathcal{Q}_U$는 일치한다.
따라서 원하는 대로 $\mathcal{F}|_U \to \mathcal{Q}_U$를 전사
$\mathcal{F} \to \mathcal{Q}$로 내릴 수 있다
(『대수공간의 성질』의 명제
\ref{spaces-properties-proposition-quasi-coherent}).
\end{proof}

\begin{lemma}
\label{moduli-lemma-quot-functorial}
$B$를 대수공간이라 하자. $\pi : X \to Y$를 $B$ 위에서 분리이고
유한 표시인 대수공간들 사이의 아핀 준유한 사상이라 하자.
$\mathcal{F}$를 준연접 $\mathcal{O}_X$-가군이라 하자. 그러면
$\pi_*$는 사상
$\Quotfunctor_{\mathcal{F}/X/B} \to \Quotfunctor_{\pi_*\mathcal{F}/Y/B}$를
유도한다.
\end{lemma}

\begin{proof}
$\mathcal{G} = \pi_*\mathcal{F}$라 두자. $\pi$가 아핀이므로 $B$ 위의
임의의 스킴 $T$에 대해 『대수공간의 코호몰로지』의 보조정리
\ref{spaces-cohomology-lemma-affine-base-change}에 의해
$\mathcal{G}_T = \pi_{T, *}\mathcal{F}_T$이다.
또한 $\pi_T$가 아핀이므로 $\pi_{T, *}$는 완전하고 몫을 몫으로 보낸다.
준연접 몫 $\mathcal{F}_T \to \mathcal{Q}$가
$\Quotfunctor_{X/B}$의 점을 정의할 필요충분조건은 $\mathcal{Q}$가
$T$ 위의 $\Cohstack_{X/B}$의 대상을 정의하는 것이다
($\mathcal{G}$와 $Y$에 대해서도 마찬가지이다).
보조정리 \ref{moduli-lemma-coherent-functorial}에서 $\pi_*$가 사상
$\Cohstack_{X/B} \to \Cohstack_{Y/B}$를 유도함을 보았으므로,
$\mathcal{F}_T \to \mathcal{Q}$가
$\Quotfunctor_{\mathcal{F}/X/B}(T)$에 속하면
$\mathcal{G}_T \to \pi_{T, *}\mathcal{Q}$는
$\Quotfunctor_{\mathcal{G}/Y/B}(T)$에 속한다.
\end{proof}

\begin{lemma}
\label{moduli-lemma-quot-open}
$B$를 대수공간이라 하자. $\pi : X \to Y$를 $B$ 위에서 분리이고
유한 표시인 대수공간들의 아핀 열린 몰입이라 하자.
$\mathcal{F}$를 준연접 $\mathcal{O}_X$-가군이라 하자. 그러면
보조정리 \ref{moduli-lemma-quot-functorial}의 사상
$\Quotfunctor_{\mathcal{F}/X/B} \to \Quotfunctor_{\pi_*\mathcal{F}/Y/B}$는
열린 몰입이다.
\end{lemma}

\begin{proof}
생략한다. 도움말: $(\pi_*\mathcal{F})_T \to \mathcal{Q}$가
$\Quotfunctor_{\pi_*\mathcal{F}/Y/B}(T)$의 원소이고 $t \in T$에 대해
$\text{Supp}(\mathcal{Q}_t) \subset |X_t|$이면, $t$의 어떤 근방에 있는
모든 $t' \in T$에 대해서도 같은 명제가 성립한다.
\end{proof}

\begin{lemma}
\label{moduli-lemma-quot-better-open}
$B$를 대수공간이라 하자. $j : X \to Y$를 $B$ 위에서 분리이고
유한 표시인 대수공간들의 열린 몰입이라 하자.
$\mathcal{G}$를 준연접 $\mathcal{O}_Y$-가군이라 하고
$\mathcal{F} = j^*\mathcal{G}$라 두자. 그러면 $B$ 위의 대수공간들의
열린 몰입
$$
\Quotfunctor_{\mathcal{F}/X/B}
\longrightarrow
\Quotfunctor_{\mathcal{G}/Y/B}
$$
이 존재한다.
\end{lemma}

\begin{proof}
$\mathcal{F}_T \to \mathcal{Q}$가
$\Quotfunctor_{\mathcal{F}/X/B}(T)$의 원소이면
$\mathcal{G}_T \to j_{T, *}\mathcal{F}_T \to j_{T, *}\mathcal{Q}$를
생각할 수 있다. 줄기를 살펴보면 이것이 전사임을 알 수 있다.
보조정리 \ref{moduli-lemma-coherent-functorial}에 의해
$j_{T, *}\mathcal{Q}$는 유한 표시이고 $B$ 위에서 평탄하며 그 지지가
$B$ 위에서 고유이다. 따라서
$\Quotfunctor_{\mathcal{G}/Y/B}$의 $T$-값 점을 얻는다.
이는 보조정리의 사상을 정의한다. 이것이 열린 몰입이라는 증명은
생략한다. 도움말: $\mathcal{G}_T \to \mathcal{Q}$가
$\Quotfunctor_{\mathcal{G}/Y/B}(T)$의 원소이고 $t \in T$에 대해
$\text{Supp}(\mathcal{Q}_t) \subset |X_t|$이면, $t$의 어떤 근방에 있는
모든 $t' \in T$에 대해서도 같은 명제가 성립한다.
\end{proof}

\begin{lemma}
\label{moduli-lemma-quot-closed}
$B$를 대수공간이라 하자. $\pi : X \to Y$를 $B$ 위에서 분리이고
유한 표시인 대수공간들의 닫힌 몰입이라 하자.
$\mathcal{F}$를 준연접 $\mathcal{O}_X$-가군이라 하자. 그러면 보조정리
\ref{moduli-lemma-quot-functorial}의 사상
$\Quotfunctor_{\mathcal{F}/X/B} \to \Quotfunctor_{\pi_*\mathcal{F}/Y/B}$는
동형이다.
\end{lemma}

\begin{proof}
$B$ 위의 모든 스킴 $T$에 대해 사상 $\pi_T : X_T \to Y_T$는 닫힌
몰입이다. 그러면 $\pi_{T, *}$는 $\QCoh(\mathcal{O}_{X_T})$와,
$X_T$의 아이디얼 층에 의해 소멸되는 준연접 가군들을 대상으로 하는
$\QCoh(\mathcal{O}_{Y_T})$의 충만한 부분범주 사이의 동치이다.
『대수공간의 사상』의 보조정리
\ref{spaces-morphisms-lemma-i-star-equivalence}를 보라.
$(\pi_*\mathcal{F})_T$의 몫은 이 아이디얼에 의해 소멸되므로 모든
$T$에 대해 사상
$\Quotfunctor_{\mathcal{F}/X/B}(T) \to
\Quotfunctor_{\pi_*\mathcal{F}/Y/B}(T)$가 전단사임을 얻는다.
\end{proof}

\begin{lemma}
\label{moduli-lemma-quot-quotient}
$X \to B$를 이 절의 도입부에서와 같이 두자.
$\mathcal{F} \to \mathcal{G}$를 준연접 $\mathcal{O}_X$-가군의 전사라
하자. 그러면 표준적인 닫힌 몰입
$\Quotfunctor_{\mathcal{G}/X/B} \to \Quotfunctor_{\mathcal{F}/X/B}$가
존재한다.
\end{lemma}

\begin{proof}
$\mathcal{K} = \Ker(\mathcal{F} \to \mathcal{G})$라 하자.
당김의 오른쪽 완전성에 의해 $B$ 위의 모든 스킴 $T$에 대해
$\mathcal{K}_T \to \mathcal{F}_T \to \mathcal{G}_T \to 0$은
완전열이다. 특히 $\mathcal{G}_T$의 몫은 $\mathcal{F}_T$의 몫을
결정하므로 함자들의 변환
$\Quotfunctor_{\mathcal{G}/X/B} \to \Quotfunctor_{\mathcal{F}/X/B}$를
얻는다. 이 변환은 『대수공간 위의 평탄성』의 보조정리
\ref{spaces-flat-lemma-F-zero-closed-proper}에 의해 닫힌 몰입이다.
구체적으로 $\Quotfunctor_{\mathcal{F}/X/B}(T)$의 원소
$\mathcal{F}_T \to \mathcal{Q}$가 주어지면, $T'/T$로의 당김이 이
변환의 상에 속할 필요충분조건은
$\mathcal{K}_{T'} \to \mathcal{Q}_{T'}$가 영인 것이다.
\end{proof}

\begin{remark}[수치 불변량]
\label{moduli-remark-quot-numerical}
$f : X \to B$와 $\mathcal{F}$를 이 절의 도입부에서와 같이 두자.
$I$를 집합이라 하고 각 $i \in I$에 대해
$E_i \in D(\mathcal{O}_X)$가 완전하다고 하자.
$P : I \to \mathbf{Z}$를 함수라 하자. 다음 사상이 있음을 상기하자.
$$
\Quotfunctor_{\mathcal{F}/X/B} \longrightarrow \Cohstack_{X/B}
$$
이 사상은 $\Quotfunctor_{\mathcal{F}/X/B}(T)$의 원소
$\mathcal{F}_T \to \mathcal{Q}$를 $T$ 위의 $\Cohstack_{X/B}$의 대상
$\mathcal{Q}$로 보낸다. 『몫』의 명제
\ref{quot-proposition-quot}의 증명을 보라. 따라서 섬유곱 도식
$$
\xymatrix{
\Quotfunctor^P_{\mathcal{F}/X/B} \ar[r] \ar[d] &
\Cohstack^P_{X/B} \ar[d] \\
\Quotfunctor_{\mathcal{F}/X/B} \ar[r] &
\Cohstack_{X/B}
}
$$
을 만들 수 있다. 이는 왼쪽 위의 대수공간을 정의하는 도식이다.
왼쪽 수직 화살표는 평탄한 닫힌 몰입이며, 예를 들어 $I$가 유한이거나,
$B$가 국소 뇌터이거나, 어떤 가역 $\mathcal{O}_X$-가군
$\mathcal{L}$에 대해 $I = \mathbf{Z}$이고
$E_i = \mathcal{L}^{\otimes i}$이면 열린 동시에 닫힌 몰입이다
(마지막 경우에는 때때로
$\Quotfunctor^{P, \mathcal{L}}_{\mathcal{F}/X/B}$라는 표기를 쓴다).
상황 \ref{moduli-situation-numerical}, 보조정리 \ref{moduli-lemma-open-P}와
\ref{moduli-lemma-finite-list-perfect-objects}, 그리고 예
\ref{moduli-example-hilbert-polynomial}을 보라.
\end{remark}

\begin{lemma}
\label{moduli-lemma-quot-tensor-invertible}
$f : X \to B$와 $\mathcal{F}$를 이 절의 도입부에서와 같이 두자.
$\mathcal{L}$을 가역 $\mathcal{O}_X$-가군이라 하자. 그러면
$\mathcal{L}$과의 텐서곱은 동형
$$
\Quotfunctor_{\mathcal{F}/X/B}
\longrightarrow
\Quotfunctor_{\mathcal{F} \otimes_{\mathcal{O}_X} \mathcal{L}/X/B}
$$
을 정의한다. 수치 다항식 $P(t)$가 주어졌을 때
$P'(t) = P(t + 1)$로 두면, 이 사상은 열린 동시에 닫힌 부분스택들의
동형
$\Quotfunctor^P_{\mathcal{F}/X/B}
\longrightarrow
\Quotfunctor^{P'}_{\mathcal{F} \otimes_{\mathcal{O}_X} \mathcal{L}/X/B}$를
유도한다.
\end{lemma}

\begin{proof}
$\mathcal{G} = \mathcal{F} \otimes_{\mathcal{O}_X} \mathcal{L}$이라 두자.
$\mathcal{G}_T = \mathcal{F}_T \otimes_{\mathcal{O}_{X_T}} \mathcal{L}_T$
임에 유의하자. $\mathcal{F}_T \to \mathcal{Q}$가
$\Quotfunctor_{\mathcal{F}/X/B}(T)$의 원소이면, 이를
$\mathcal{G}_T \to
\mathcal{Q} \otimes_{\mathcal{O}_{X_T}} \mathcal{L}_T$라는
$\Quotfunctor_{\mathcal{F} \otimes_{\mathcal{O}_X} \mathcal{L}/X/B}(T)$의
원소로 보낸다. 이는 당김과 호환되므로 원하는 함자들의 변환을 정의한다.
명백한 역변환이 있으므로 동형이다. 마지막 명제의 증명은 생략한다.
\end{proof}

\begin{lemma}
\label{moduli-lemma-quot-power-invertible}
$f : X \to B$와 $\mathcal{F}$를 이 절의 도입부에서와 같이 두자.
$\mathcal{L}$을 가역 $\mathcal{O}_X$-가군이라 하자. 그러면
$$
\Quotfunctor^{P, \mathcal{L}}_{\mathcal{F}/X/B} =
\Quotfunctor^{P', \mathcal{L}^{\otimes n}}_{\mathcal{F}/X/B}
$$
이며, 여기서 $P'(t) = P(nt)$이다.
\end{lemma}

\begin{proof}
모든 정의를 풀어 쓰면 즉시 따른다.
\end{proof}




\section{Quot의 유계성}
\label{moduli-section-quot-bounded}

\noindent
고전적인 상황과 달리 Quot 함자가 대수공간이라는 것은 이미 알지만,
그것이 유한형 대수공간으로 표현되는 경우가 있는지는 아직 모른다.

\begin{lemma}
\label{moduli-lemma-quot-Pn}
$n \geq 0$, $r \geq 1$, $P \in \mathbf{Q}[t]$라 하자.
힐베르트 다항식이 $P$인
$\mathcal{O}_{\mathbf{P}^n_\mathbf{Z}}^{\oplus r}$의 몫을 매개화하는
대수공간
$$
X = \Quotfunctor^P_{\mathcal{O}^{\oplus r}_{\mathbf{P}^n_\mathbf{Z}}/
\mathbf{P}^n_\mathbf{Z}/\mathbf{Z}}
$$
는 $\Spec(\mathbf{Z})$ 위에서 고유이다.
\end{lemma}

\begin{proof}
$X \to \Spec(\mathbf{Z})$가 분리이고 국소적으로 유한 표시임은 이미
안다(보조정리 \ref{moduli-lemma-quot-s-lfp}). 또한
$X \to \Spec(\mathbf{Z})$가 값매김 판정법의 존재 부분을 만족함도
안다. 보조정리 \ref{moduli-lemma-quot-existence-part}를 보라.
고유성에 대한 값매김 판정법에 의해 우리의 Quot 공간이 준콤팩트임을
보이면 충분하다. 『대수공간의 사상』의 보조정리
\ref{spaces-morphisms-lemma-characterize-proper}를 보라.
따라서 준콤팩트 스킴 $T$와 전사 사상 $T \to X$를 찾으면 충분하다.
『다양체』의 보조정리 \ref{varieties-lemma-bound-quotients-free}에서 찾은
정수를 $m$이라 하자. 다음과 같이 두자.
$$
N = r{m + n \choose n} - P(m)
$$
$\mathbf{P}^n_\mathbf{Z} = \text{Proj}(\mathbf{Z}[T_0, \ldots, T_n])$를
$\mathbf{P}^n$으로 쓰고, 장식이 없는 곱은 $\Spec(\mathbf{Z})$ 위의
곱을 뜻하게 하겠다. 증명의 발상은 아핀 스킴 $T$ 위에서
``보편적인'' 사상
$$
\Psi :
\mathcal{O}_{T \times \mathbf{P}^n}(-m)^{\oplus N}
\longrightarrow
\mathcal{O}_{T \times \mathbf{P}^n}^{\oplus r}
$$
을 구성하고, $X$의 모든 점이 $T$의 어떤 점에서 이 사상의 여핵에
대응함을 보이는 것이다.

\medskip\noindent
$T$와 $\Psi$의 정의. $T = \Spec(A)$로 두는데, 여기서
$$
A = \mathbf{Z}[a_{i, j, E}]
$$
이고 $i \in \{1, \ldots, r\}$, $j \in \{1, \ldots, N\}$이며
$E = (e_0, \ldots, e_n)$는 전차수가
$|E| = \sum_{k = 0, \ldots n} e_k = m$인 다중지표를 달린다.
그런 다음 $\Psi$를 그 $(i, j)$ 행렬 성분이 사상
$$
\sum\nolimits_{E = (e_0, \ldots, e_n)}
a_{i, j, E} T_0^{e_0} \ldots T_n^{e_n} :
\mathcal{O}_{T \times \mathbf{P}^n}(-m)
\longrightarrow
\mathcal{O}_{T \times \mathbf{P}^n}
$$
이 되도록 정의한다. 여기서 합은 위와 같은 $E$에 걸쳐 취한다
(물론 $i$와 $j$는 고정되어 있다).

\medskip\noindent
$T \times \mathbf{P}^n$ 위의 몫
$\mathcal{Q} = \Coker(\Psi)$를 생각하자. 『사상 더 보기』의 보조정리
\ref{more-morphisms-lemma-generic-flatness-stratification}에 의해
$t \geq 0$와 닫힌 부분스킴들
$$
T = T_0 \supset T_1 \supset \ldots \supset T_t = \emptyset
$$
이 존재하여, $\mathcal{Q}$를
$(T_p \setminus T_{p + 1}) \times \mathbf{P}^n$으로 당긴
$\mathcal{Q}_p$는 $T_p \setminus T_{p + 1}$ 위에서 평탄하다.
$\mathcal{Q} = \Coker(\Psi)$를 정의하는 완전열을 당기면 완전열
$$
\mathcal{O}_{(T_p \setminus T_{p + 1}) \times \mathbf{P}^n}(-m)^{\oplus N}
\to
\mathcal{O}_{(T_p \setminus T_{p + 1}) \times \mathbf{P}^n}^{\oplus r}
\to
\mathcal{Q}_p
\to
0
$$
을 얻는다. 따라서 사상
$$
\coprod (T_p \setminus T_{p + 1})
\longrightarrow
\Quotfunctor_{\mathcal{O}^{\oplus r}/\mathbf{P}/\mathbf{Z}}
\supset
\Quotfunctor^P_{\mathcal{O}^{\oplus r}/\mathbf{P}/\mathbf{Z}} = X
$$
을 얻는다. 왼쪽은 뇌터 스킴이고 오른쪽의 포함은 열려 있으므로,
$X$의 임의의 점이 이 사상의 상에 속함을 보이면 충분하다.

\medskip\noindent
$k$를 체라 하고 $x \in X(k)$라 하자. 그러면 $x$는 힐베르트 다항식이
$P$인 연접 $\mathcal{O}_{\mathbf{P}^n_k}$-가군 $\mathcal{F}$로의 전사
$\mathcal{O}_{\mathbf{P}^n_k}^{\oplus r} \to \mathcal{F}$에 대응한다.
짧은 완전열
$$
0 \to \mathcal{K} \to
\mathcal{O}_{\mathbf{P}^n_k}^{\oplus r} \to
\mathcal{F} \to 0
$$
을 생각하자. 『다양체』의 보조정리
\ref{varieties-lemma-bound-quotients-free}와 $m$의 선택에 의해
$\mathcal{K}$는 $m$-정칙이다. 『다양체』의 보조정리
\ref{varieties-lemma-m-regular-globally-generated}에 의해
$\mathcal{K}(m)$은 대역적으로 생성된다.
『다양체』의 보조정리 \ref{varieties-lemma-m-regular-up}과
$m$-정칙성의 정의에 의해 $i > 0$이면
$H^i(\mathbf{P}^n_k, \mathcal{K}(m)) = 0$이다. 따라서 $N$의 선택에 의해
$$
\dim_k H^0(\mathbf{P}^n_k, \mathcal{K}(m)) =
\chi(\mathcal{K}(m)) =
\chi(\mathcal{O}_{\mathbf{P}^n_k}(m)^{\oplus r}) -
\chi(\mathcal{F}(m)) = N
$$
임을 알 수 있다. 이로부터 전사
$$
\mathcal{O}_{\mathbf{P}^n_k}^{\oplus N}
\longrightarrow
\mathcal{K}(m)
$$
를 얻는다. 다시 아래로 꼬고 위의 짧은 완전열을 사용하면
$\mathcal{F}$가 사상
$$
\Psi_x :
\mathcal{O}_{\mathbf{P}^n_k}(-m)^{\oplus N}
\to
\mathcal{O}_{\mathbf{P}^n_k}^{\oplus r}
$$
의 여핵임을 알 수 있다. 대응하는 사상
$t = \Spec(\tau) : \Spec(k) \to T$에 의한 $\Psi$의 밑변환이
$\Psi_x$가 되게 하는 환 준동형 $\tau : A \to k$가 유일하게 존재한다.
이는 $\Psi_x$를 정의하는 $N \times r$ 행렬의 성분들이
$T_0, \ldots, T_n$에 관한 차수 $m$의 동차다항식
$\sum \lambda_{i, j, E} T_0^{e_0} \ldots T_n^{e_n}$이고 계수들이
$\lambda_{i, j, E} \in k$이므로
$\tau(a_{i, j, E}) = \lambda_{i, j, E}$로 둘 수 있기 때문이다.
그러면 어떤 $p$에 대해 $t \in T_p \setminus T_{p + 1}$이고,
위의 사상 아래에서 $t$의 상은 원하는 대로 $x$이다.
\end{proof}

\begin{lemma}
\label{moduli-lemma-quot-Pn-over-base}
$B$를 대수공간이라 하고 $X = B \times \mathbf{P}^n_\mathbf{Z}$라 하자.
$\mathcal{L}$을 $\mathcal{O}_{\mathbf{P}^n}(1)$의 $X$로의 당김이라 하자.
$\mathcal{F}$를 유한 표시 $\mathcal{O}_X$-가군이라 하자.
$\mathcal{L}$에 관해 힐베르트 다항식이 $P$인 $\mathcal{F}$의 몫을
매개화하는 대수공간 $\Quotfunctor^P_{\mathcal{F}/X/B}$는
$B$ 위에서 고유이다.
\end{lemma}

\begin{proof}
이 문제는 $B$ 위에서 에탈 국소적이다. 『대수공간의 사상』의 보조정리
\ref{spaces-morphisms-lemma-proper-local}을 보라. 따라서 $B$가 아핀
스킴이라고 가정할 수 있다. 이 경우 $\mathcal{L}$은 $X$ 위의 풍부한
가역 가군이다(『구성』의 보조정리
\ref{constructions-lemma-ample-on-proj}와 『성질』의 정의
\ref{properties-definition-ample}에 나오는 풍부한 가역 가군의 정의에
의한다). 따라서 『성질』의 명제
\ref{properties-proposition-characterize-ample}에 의해
$r' \geq 0$와 $r \geq 0$ 및 전사
$$
\mathcal{O}_X^{\oplus r} \longrightarrow
\mathcal{F} \otimes_{\mathcal{O}_X} \mathcal{L}^{\otimes r'}
$$
를 찾을 수 있다. 보조정리 \ref{moduli-lemma-quot-tensor-invertible}에 의해
$\mathcal{F}$를
$\mathcal{F} \otimes_{\mathcal{O}_X} \mathcal{L}^{\otimes r'}$로,
$P(t)$를 $P(t + r')$로 바꿀 수 있다.
보조정리 \ref{moduli-lemma-quot-quotient}에 의해 닫힌 몰입
$$
\Quotfunctor^P_{\mathcal{F}/X/B}
\longrightarrow
\Quotfunctor^P_{\mathcal{O}_X^{\oplus r}/X/B}
$$
을 얻는다. 보조정리 \ref{moduli-lemma-quot-Pn}에서
$\Quotfunctor^P_{\mathcal{O}_X^{\oplus r}/X/B} \to B$가 고유임을
보였으므로 결론을 얻는다.
\end{proof}

\begin{lemma}
\label{moduli-lemma-quot-proper-over-base}
$f : X \to B$를 대수공간의 고유 유한 표시 사상이라 하자.
$\mathcal{F}$를 유한 표시 $\mathcal{O}_X$-가군이라 하자.
$\mathcal{L}$을 $X/B$ 위에서 풍부한 가역 $\mathcal{O}_X$-가군이라 하자.
『대수공간 위의 인자』의 정의
\ref{spaces-divisors-definition-relatively-ample}를 보라.
$\mathcal{L}$에 관해 힐베르트 다항식이 $P$인 $\mathcal{F}$의 몫을
매개화하는 대수공간 $\Quotfunctor^P_{\mathcal{F}/X/B}$는
$B$ 위에서 고유이다.
\end{lemma}

\begin{proof}
이 문제는 $B$ 위에서 에탈 국소적이다. 『대수공간의 사상』의 보조정리
\ref{spaces-morphisms-lemma-proper-local}을 보라. 따라서 $B$가 아핀
스킴이라고 가정할 수 있다. 그러면 어떤 $d \geq 1$에 대해
$i^*\mathcal{O}_{\mathbf{P}^n_B}(1) \cong \mathcal{L}^{\otimes d}$를
만족하는 닫힌 몰입 $i : X \to \mathbf{P}^n_B$를 찾을 수 있다.
『사상』의 보조정리
\ref{morphisms-lemma-quasi-projective-finite-type-over-S}를 보라.
$\mathcal{L}$을 $\mathcal{L}^{\otimes d}$로, 수치 다항식 $P(t)$를
$P(dt)$로 바꾸어도 $\Quotfunctor^P_{\mathcal{F}/X/B}$는 변하지 않는다.
일부 세부사항은 생략한다. 따라서
$\mathcal{L} = i^*\mathcal{O}_{\mathbf{P}^n_B}(1)$이라 가정할 수 있다.
그러면 보조정리 \ref{moduli-lemma-quot-closed}의 동형
$\Quotfunctor_{\mathcal{F}/X/B} \to
\Quotfunctor_{i_*\mathcal{F}/\mathbf{P}^n_B/B}$는 동형
$\Quotfunctor^P_{\mathcal{F}/X/B} \cong
\Quotfunctor^P_{i_*\mathcal{F}/\mathbf{P}^n_B/B}$를 유도한다.
보조정리 \ref{moduli-lemma-quot-Pn-over-base}에 의해
$\Quotfunctor^P_{i_*\mathcal{F}/\mathbf{P}^n_B/B}$가 $B$ 위에서
고유이므로 결론을 얻는다.
\end{proof}

\begin{lemma}
\label{moduli-lemma-quot-qc-over-base}
$f : X \to B$를 대수공간의 분리 유한 표시 사상이라 하자.
$\mathcal{F}$를 유한 표시 $\mathcal{O}_X$-가군이라 하자.
$\mathcal{L}$을 $X/B$ 위에서 풍부한 가역 $\mathcal{O}_X$-가군이라 하자.
『대수공간 위의 인자』의 정의
\ref{spaces-divisors-definition-relatively-ample}를 보라.
$\mathcal{L}$에 관해 힐베르트 다항식이 $P$인 $\mathcal{F}$의 몫을
매개화하는 대수공간 $\Quotfunctor^P_{\mathcal{F}/X/B}$는
$B$ 위에서 분리이고 유한 표시이다.
\end{lemma}

\begin{proof}
$\Quotfunctor_{\mathcal{F}/X/B} \to B$가 분리이고 국소적으로 유한
표시임은 이미 보았다. 보조정리 \ref{moduli-lemma-quot-s-lfp}를 보라.
따라서 주석 \ref{moduli-remark-quot-numerical}의 열린 부분공간
$\Quotfunctor^P_{\mathcal{F}/X/B}$가 $B$ 위에서 준콤팩트임을 보이면
충분하다.

\medskip\noindent
이 문제는 $B$ 위에서 에탈 국소적이다
(『대수공간의 사상』의 보조정리
\ref{spaces-morphisms-lemma-quasi-compact-local}). 따라서 $B$가
아핀이라고 가정할 수 있다.

\medskip\noindent
$B = \Spec(\Lambda)$라 하자. $\Lambda = \colim \Lambda_i$를 그 유한형
$\mathbf{Z}$-부분대수들의 쌍대극한으로 쓰자. 그러면 어떤 $i$에 대해
$B_i = \Spec(\Lambda_i)$ 위에서 보조정리의 조건을 만족하는 계
$X_i, \mathcal{F}_i, \mathcal{L}_i$를 찾아, 이를 $B$로 밑변환하면
$X, \mathcal{F}, \mathcal{L}$이 되게 할 수 있다. 이는
『대수공간의 극한』의 보조정리
\ref{spaces-limits-lemma-descend-finite-presentation}($X_i$를 찾는 데),
\ref{spaces-limits-lemma-descend-modules-finite-presentation}
($\mathcal{F}_i$를 찾는 데),
\ref{spaces-limits-lemma-descend-invertible-modules}
($\mathcal{L}_i$를 찾는 데), 그리고
\ref{spaces-limits-lemma-descend-separated}($X_i$를 분리로 만드는 데)에서
따른다. 다음 등식과 $\Quotfunctor^P_{\mathcal{F}/X/B}$에 대한 같은
등식에 의해 다음 문단에서 다루는 경우로 귀착된다.
$$
\Quotfunctor_{\mathcal{F}/X/B} = B \times_{B_i}
\Quotfunctor_{\mathcal{F}_i/X_i/B_i}
$$

\medskip\noindent
$B$가 뇌터 아핀 스킴이라고 하자. 보조정리
\ref{moduli-lemma-quot-power-invertible}에 따라 $\mathcal{L}$을 양의 거듭제곱으로
바꿀 수 있다. 따라서 $i^*\mathcal{O}_{\mathbf{P}^n}(1) = \mathcal{L}$을
만족하는 몰입 $i : X \to \mathbf{P}^n_B$가 존재한다고 가정할 수 있다.
『사상』의 보조정리 \ref{morphisms-lemma-quasi-compact-immersion}에 의해
$i$가 열린 몰입 $j : X \to X'$를 통해 분해되는 닫힌 부분스킴
$X' \subset \mathbf{P}^n_B$가 존재한다. 『성질』의 보조정리
\ref{properties-lemma-lift-finite-presentation}에 의해
$j^*\mathcal{G} = \mathcal{F}$를 만족하는 유한 표시
$\mathcal{O}_{X'}$-가군 $\mathcal{G}$가 존재한다.
따라서 보조정리 \ref{moduli-lemma-quot-better-open}에 의해 열린 몰입
$$
\Quotfunctor_{\mathcal{F}/X/B}
\longrightarrow
\Quotfunctor_{\mathcal{G}/X'/B}
$$
을 얻는다. 명백히 이 열린 몰입은
$\Quotfunctor^P_{\mathcal{F}/X/B}$를
$\Quotfunctor^P_{\mathcal{G}/X'/B}$ 안으로 보낸다.
이제 보조정리 \ref{moduli-lemma-quot-proper-over-base}에 의해
$\Quotfunctor^P_{\mathcal{G}/X'/B}$는 $B$ 위에서 고유이다.
따라서 이 공간은 뇌터이고, 뇌터 대수공간의 임의의 열린 부분공간은
준콤팩트이므로 결론을 얻는다.
\end{proof}





\section{힐베르트 함자의 성질}
\label{moduli-section-hilb}

\noindent
$f : X \to B$를 분리이고 유한 표시인 대수공간의 사상이라 하자.
그러면 $\Hilbfunctor_{X/B}$는 $B$ 위에서 국소적으로 유한 표시인
대수공간이다. 『몫』의 명제 \ref{quot-proposition-hilb}를 보라.

\begin{lemma}
\label{moduli-lemma-hilb-diagonal-closed}
$\Hilbfunctor_{X/B} \to B$의 대각사상은 유한 표시인 닫힌 몰입이다.
\end{lemma}

\begin{proof}
『몫』의 보조정리 \ref{quot-lemma-hilb-is-quot}에서
$\Hilbfunctor_{X/B} = \Quotfunctor_{\mathcal{O}_X/X/B}$임을 보았다.
따라서 이는 보조정리 \ref{moduli-lemma-quot-diagonal-closed}에서 따른다.
\end{proof}

\begin{lemma}
\label{moduli-lemma-hilb-s-lfp}
사상 $\Hilbfunctor_{X/B} \to B$는 분리이고 국소적으로 유한 표시이다.
\end{lemma}

\begin{proof}
$\Hilbfunctor_{X/B} \to B$가 분리임을 확인하려면 그 대각사상이
닫힌 몰입임을 보여야 한다. 이는 보조정리
\ref{moduli-lemma-hilb-diagonal-closed}에 의해 참이다.
두 번째 명제는 『몫』의 명제 \ref{quot-proposition-hilb}의 일부이다.
\end{proof}

\begin{lemma}
\label{moduli-lemma-hilb-existence-part}
$X \to B$가 유한 표시일 뿐 아니라 고유라고 하자. 그러면
$\Hilbfunctor_{X/B} \to B$는 값매김 판정법의 존재 부분을 만족한다
(『대수공간의 사상』의 정의
\ref{spaces-morphisms-definition-valuative-criterion}).
\end{lemma}

\begin{proof}
『몫』의 보조정리 \ref{quot-lemma-hilb-is-quot}에서
$\Hilbfunctor_{X/B} = \Quotfunctor_{\mathcal{O}_X/X/B}$임을 보았다.
따라서 이는 보조정리 \ref{moduli-lemma-quot-existence-part}에서 따른다.
\end{proof}

\begin{lemma}
\label{moduli-lemma-hilb-open}
$B$를 대수공간이라 하자. $\pi : X \to Y$를 $B$ 위에서 분리이고
유한 표시인 대수공간들의 열린 몰입이라 하자. 그러면 $\pi$는 열린 몰입
$\Hilbfunctor_{X/B} \to \Hilbfunctor_{Y/B}$를 유도한다.
\end{lemma}

\begin{proof}
생략한다. 도움말: $Z \subset X_T$가 $T$ 위에서 고유인 닫힌
부분스킴이면 $Z$는 $Y_T$에서도 닫혀 있다. 따라서 변환
$\Hilbfunctor_{X/B} \to \Hilbfunctor_{Y/B}$를 얻는다.
$Z \subset Y_T$가 $\Hilbfunctor_{Y/B}(T)$의 원소이고 $t \in T$에 대해
$|Z_t| \subset |X_t|$이면, $t$의 어떤 근방에 있는 모든 $t' \in T$에
대해서도 같은 명제가 성립한다.
\end{proof}

\begin{lemma}
\label{moduli-lemma-hilb-closed}
$B$를 대수공간이라 하자. $\pi : X \to Y$를 $B$ 위에서 분리이고
유한 표시인 대수공간들의 닫힌 몰입이라 하자. 그러면 $\pi$는 닫힌 몰입
$\Hilbfunctor_{X/B} \to \Hilbfunctor_{Y/B}$를 유도한다.
\end{lemma}

\begin{proof}
$\pi$가 닫힌 몰입이므로, 닫힌 부분스킴 $Z \subset X_T$가 주어지면
$Z$를 $X_T$의 닫힌 부분스킴으로 볼 수 있음은 즉시 알 수 있다.
따라서 변환 $\Hilbfunctor_{X/B} \to \Hilbfunctor_{Y/B}$를 얻는다.
이 변환이 단사 사상임도 즉시 알 수 있다. 이것이 닫힌 몰입임을
증명하려면 사상 $\mathcal{O}_Y \to \mathcal{O}_X$에 보조정리
\ref{moduli-lemma-quot-quotient}을 적용하고, 『몫』의 보조정리
\ref{quot-lemma-hilb-is-quot}의 식별
$\Hilbfunctor_{X/B} = \Quotfunctor_{\mathcal{O}_X/X/B}$,
$\Hilbfunctor_{Y/B} = \Quotfunctor_{\mathcal{O}_Y/Y/B}$를 사용하면 된다.
\end{proof}

\begin{remark}[수치 불변량]
\label{moduli-remark-hilb-numerical}
$f : X \to B$를 이 절의 도입부에서와 같이 두자. $I$를 집합이라 하고
각 $i \in I$에 대해 $E_i \in D(\mathcal{O}_X)$가 완전하다고 하자.
$P : I \to \mathbf{Z}$를 함수라 하자. 『몫』의 보조정리
\ref{quot-lemma-hilb-is-quot}에서
$\Hilbfunctor_{X/B} = \Quotfunctor_{\mathcal{O}_X/X/B}$임을 상기하자.
따라서
$$
\Hilbfunctor^P_{X/B} = \Quotfunctor^P_{\mathcal{O}_X/X/B}
$$
로 정의할 수 있으며, 여기서
$\Quotfunctor^P_{\mathcal{O}_X/X/B}$는 주석
\ref{moduli-remark-quot-numerical}에서와 같다. 사상
$$
\Hilbfunctor^P_{X/B} \longrightarrow \Hilbfunctor_{X/B}
$$
은 평탄한 닫힌 몰입이며, 예를 들어 $I$가 유한이거나 $B$가 국소
뇌터이거나 어떤 가역 $\mathcal{O}_X$-가군 $\mathcal{L}$에 대해
$I = \mathbf{Z}$이고 $E_i = \mathcal{L}^{\otimes i}$이면 열린 동시에
닫힌 몰입이다. 마지막 경우에는 때때로
$\Hilbfunctor^{P, \mathcal{L}}_{X/B}$라는 표기를 쓴다.
\end{remark}

\begin{lemma}
\label{moduli-lemma-hilb-proper-over-base}
$f : X \to B$를 대수공간의 고유 유한 표시 사상이라 하자.
$\mathcal{L}$을 $X/B$ 위에서 풍부한 가역 $\mathcal{O}_X$-가군이라 하자.
『대수공간 위의 인자』의 정의
\ref{spaces-divisors-definition-relatively-ample}를 보라.
$\mathcal{L}$에 관해 힐베르트 다항식이 $P$인 닫힌 부분스킴을
매개화하는 대수공간 $\Hilbfunctor^P_{X/B}$는 $B$ 위에서 고유이다.
\end{lemma}

\begin{proof}
『몫』의 보조정리 \ref{quot-lemma-hilb-is-quot}에서
$\Hilbfunctor_{X/B} = \Quotfunctor_{\mathcal{O}_X/X/B}$임을 상기하자.
따라서 이 보조정리는 보조정리
\ref{moduli-lemma-quot-proper-over-base}의 직접적인 귀결이다.
\end{proof}

\begin{lemma}
\label{moduli-lemma-hilb-qc-over-base}
$f : X \to B$를 대수공간의 분리 유한 표시 사상이라 하자.
$\mathcal{L}$을 $X/B$ 위에서 풍부한 가역 $\mathcal{O}_X$-가군이라 하자.
『대수공간 위의 인자』의 정의
\ref{spaces-divisors-definition-relatively-ample}를 보라.
$\mathcal{L}$에 관해 힐베르트 다항식이 $P$인 닫힌 부분스킴을
매개화하는 대수공간 $\Hilbfunctor^P_{X/B}$는 $B$ 위에서 분리이고
유한 표시이다.
\end{lemma}

\begin{proof}
『몫』의 보조정리 \ref{quot-lemma-hilb-is-quot}에서
$\Hilbfunctor_{X/B} = \Quotfunctor_{\mathcal{O}_X/X/B}$임을 상기하자.
따라서 이 보조정리는 보조정리
\ref{moduli-lemma-quot-qc-over-base}의 직접적인 귀결이다.
\end{proof}









\section{피카르 스택의 성질}
\label{moduli-section-picard-stack}

\noindent
$f : X \to B$를 평탄하고 고유이며 유한 표시인 대수공간의 사상이라
하자. 그러면 $X/B$ 위의 가역층을 매개화하는 스택
$\Picardstack_{X/B}$는 대수적이다. 『몫』의 명제
\ref{quot-proposition-pic}을 보라.

\begin{lemma}
\label{moduli-lemma-pic-diagonal-affine-fp}
$B$ 위에서 $\Picardstack_{X/B}$의 대각사상은 아핀이며 유한 표시이다.
\end{lemma}

\begin{proof}
『몫』의 보조정리 \ref{quot-lemma-picard-stack-open-in-coh}에서
$\Picardstack_{X/B}$가 $\Cohstack_{X/B}$의 열린 부분스택임을 보았다.
따라서 이는 보조정리 \ref{moduli-lemma-coherent-diagonal-affine-fp}에서 따른다.
\end{proof}

\begin{lemma}
\label{moduli-lemma-pic-qs-lfp}
사상 $\Picardstack_{X/B} \to B$는 준분리이고 국소적으로 유한 표시이다.
\end{lemma}

\begin{proof}
『몫』의 보조정리 \ref{quot-lemma-picard-stack-open-in-coh}에서
$\Picardstack_{X/B}$가 $\Cohstack_{X/B}$의 열린 부분스택임을 보았다.
따라서 이는 보조정리 \ref{moduli-lemma-coherent-qs-lfp}에서 따른다.
\end{proof}

\begin{lemma}
\label{moduli-lemma-pic-existence-part}
$X \to B$가 고유일 뿐 아니라 매끄럽다고 하자. 그러면
$\Picardstack_{X/B} \to B$는 값매김 판정법의 존재 부분을 만족한다
(『스택의 사상』의 정의 \ref{stacks-morphisms-definition-existence}).
\end{lemma}

\begin{proof}
밑변환을 취하면 곧바로 다음 문제로 귀착된다. 분수체가 $K$인
값매김환 $R$, $R$ 위에서 고유하고 매끄러운 대수공간 $X$, 그리고 가역
$\mathcal{O}_{X_K}$-가군 $\mathcal{L}_K$가 주어졌을 때, 일반 섬유가
$\mathcal{L}_K$인 가역 $\mathcal{O}_X$-가군 $\mathcal{L}$이 존재함을
보여라. $X_K$는 뇌터이고 분리이며 정칙임에 유의하자
(『대수공간의 사상』의 보조정리
\ref{spaces-morphisms-lemma-finite-presentation-noetherian}과
『체 위의 대수공간』의 보조정리
\ref{spaces-over-fields-lemma-smooth-regular}를 사용하라).
따라서 $X_K$ 안의 두 유효 카르티에 인자 $D_K, D'_K$에 대해
$\mathcal{L}_K$를 피카르 군에서 $\mathcal{O}_{X_K}(D_K)$와
$\mathcal{O}_{X_K}(D'_K)$의 차로 쓸 수 있다. 『대수공간 위의 인자』의
보조정리
\ref{spaces-divisors-lemma-Noetherian-regular-separated-pic-effective-Cartier}를
보라. 마지막으로 $D_K$와 $D'_K$가 유효 카르티에 인자
$D, D' \subset X$의 제한임을 안다. 『대수공간 위의 인자』의 보조정리
\ref{spaces-divisors-lemma-smooth-over-valuation-ring-effective-Cartier}를
보라.
\end{proof}

\begin{lemma}
\label{moduli-lemma-pic-inertia}
$B$ 위의 모든 스킴 $T$에 대해
$f_{T, *}\mathcal{O}_{X_T} \cong \mathcal{O}_T$라고 하자. 그러면
$\Picardstack_{X/B}$의 관성 스택은
$\mathbf{G}_m \times \Picardstack_{X/B}$와 같다.
\end{lemma}

\begin{proof}
이는 『스택의 예』의 예
\ref{examples-stacks-example-inertia-stack-of-picard}에서 설명한다.
\end{proof}

\begin{lemma}
\label{moduli-lemma-pic-curves-smooth}
이 절의 다른 가정들에 더하여 $f : X \to B$의 상대 차원이
$\leq 1$이라고 하자. 그러면 $\Picardstack_{X/B} \to B$는 매끄럽다.
\end{lemma}

\begin{proof}
$\Picardstack_{X/B} \to B$가 국소적으로 유한 표시임은 이미 안다.
보조정리 \ref{moduli-lemma-pic-qs-lfp}를 보라. 따라서 『스택의 사상 더 보기』의
보조정리 \ref{stacks-more-morphisms-lemma-smooth-formally-smooth}에 의해
$\Picardstack_{X/B} \to B$가 형식적으로 매끄러움을 보이면 충분하다.
밑변환을 취하면 곧바로 다음 문제로 귀착된다. 아핀 스킴의 일차 두꺼워짐
$T \subset T'$, 고유하고 평탄하며 유한 표시이고 상대 차원이
$\leq 1$인 사상 $X' \to T'$, 그리고
$X = T \times_{T'} X'$ 위의 가역 $\mathcal{O}_X$-가군
$\mathcal{L}$이 주어졌을 때, $X$로 제한하면 $\mathcal{L}$이 되는 가역
$\mathcal{O}_{X'}$-가군 $\mathcal{L}'$이 존재함을 증명하라.
$T \subset T'$가 일차 두꺼워짐이므로 $X \subset X'$도 그러하다.
『대수공간의 사상 더 보기』의 보조정리
\ref{spaces-more-morphisms-lemma-base-change-thickening}을 보라.
『대수공간의 사상 더 보기』의 보조정리
\ref{spaces-more-morphisms-lemma-picard-group-first-order-thickening}에 의해
$X$를 $X'$에서 잘라내는 준연접 아이디얼을 $\mathcal{I}$라 할 때
$H^2(X, \mathcal{I}) = 0$임을 보이면 충분하다.
구조 사상을 $f : X \to T$로 표기하자. 『대수공간의 코호몰로지』의
보조정리
\ref{spaces-cohomology-lemma-higher-direct-images-zero-above-dimension-fibre}에
의해 $p > 1$이면 $R^pf_*\mathcal{I} = 0$이다. 따라서
『대수공간의 코호몰로지』의 보조정리
\ref{spaces-cohomology-lemma-quasi-coherence-higher-direct-images-application}에
의해 원하는 소멸을 얻는다(여기서 마침내 $T$가 아핀이라는 사실을 쓴다).
\end{proof}



\section{피카르 함자의 성질}
\label{moduli-section-picard-functor}

\noindent
$f : X \to B$를 평탄하고 고유이며 유한 표시인 대수공간의 사상이라
하고, 더 나아가 모든 $T/B$에 대해 표준 사상
$$
\mathcal{O}_T \longrightarrow f_{T, *}\mathcal{O}_{X_T}
$$
이 동형이라고 하자. 그러면 피카르 함자 $\Picardfunctor_{X/B}$는
대수공간이다. 『몫』의 명제 \ref{quot-proposition-pic-functor}를 보라.
피카르 스택과는 밀접한 관계가 있다.

\begin{lemma}
\label{moduli-lemma-pic-gerbe-over-pic-functor}
사상 $\Picardstack_{X/B} \to \Picardfunctor_{X/B}$는 피카르 스택을
피카르 함자 위의 gerbe로 만든다.
\end{lemma}

\begin{proof}
$\Picardstack_{X/B} \to \Picardfunctor_{X/B}$가 gerbe라는 정의는
『스택의 사상』의 정의 \ref{stacks-morphisms-definition-gerbe}에 주어져
있고, 이는 다시 『스택』의 정의
\ref{stacks-definition-gerbe-over-stack-in-groupoids}를 참조한다.
이를 증명하기 위해 『스택』의 보조정리 \ref{stacks-lemma-when-gerbe}의
조건 (2)(a)와 (2)(b)를 확인하겠다. 이는 『몫』의 보조정리
\ref{quot-lemma-pic-over-pic}에서 즉시 따른다. 자세히 설명하자.

\medskip\noindent
조건 (2)(a). $B$ 위의 어떤 스킴 $U$에 대해
$\xi \in \Picardfunctor_{X/B}(U)$라고 하자.
$\Picardfunctor_{X/B}$는 $B$ 위의 스킴에 주어진 규칙
$T \mapsto \Pic(X_T)$의 fppf 층화이므로
(『몫』의 상황 \ref{quot-situation-pic}), fppf 덮개
$\{U_i \to U\}$가 존재하여 $\xi|_{U_i}$는 $X_{U_i}$ 위의 어떤 가역
가군 $\mathcal{L}_i$에 대응한다. 그러면
$(U_i \to B, \mathcal{L}_i)$는 $U_i$ 위의
$\Picardstack_{X/B}$의 대상이고 $\xi|_{U_i}$로 사상된다.

\medskip\noindent
조건 (2)(b). $U$가 $B$ 위의 스킴이고 $\mathcal{L}, \mathcal{N}$이
$X_U$ 위의 가역 가군으로서 $\Picardfunctor_{X/B}(U)$의 같은 원소로
사상된다고 하자. 그러면 fppf 덮개 $\{U_i \to U\}$가 존재하여
$\mathcal{L}|_{X_{U_i}}$는 $\mathcal{N}|_{X_{U_i}}$와 동형이다.
따라서 원하는 대로
$(U \to B, \mathcal{L})|_{U_i} \to (U \to B, \mathcal{N})|_{U_i}$인
동형들을 찾는다.
\end{proof}

\begin{lemma}
\label{moduli-lemma-pic-functor-diagonal-qc-immersion}
$B$ 위에서 $\Picardfunctor_{X/B}$의 대각사상은 준콤팩트 몰입이다.
\end{lemma}

\begin{proof}
『몫』의 보조정리 \ref{quot-lemma-diagonal-pic}에 의해 대각사상은
몰입이다. 마무리를 위해 대각사상이 준콤팩트임을 보이자.
보조정리 \ref{moduli-lemma-pic-diagonal-affine-fp}에 의해
$\Picardstack_{X/B}$의 대각사상은 준콤팩트이고, 보조정리
\ref{moduli-lemma-pic-gerbe-over-pic-functor}에 의해
$\Picardstack_{X/B}$는 $\Picardfunctor_{X/B}$ 위의 gerbe이다.
『스택의 사상』의 보조정리
\ref{stacks-morphisms-lemma-gerbe-diagonal-quasi-compact}로 결론을 얻는다.
\end{proof}

\begin{lemma}
\label{moduli-lemma-pic-functor-qs-lfp}
사상 $\Picardfunctor_{X/B} \to B$는 준분리이고 국소적으로 유한 표시이다.
\end{lemma}

\begin{proof}
$\Picardfunctor_{X/B} \to B$가 준분리임을 확인하려면 그 대각사상이
준콤팩트임을 보여야 한다. 이는 보조정리
\ref{moduli-lemma-pic-functor-diagonal-qc-immersion}에서 즉시 따른다.
사상 $\Picardstack_{X/B} \to \Picardfunctor_{X/B}$는 전사이고 평탄하며
국소적으로 유한 표시이므로(보조정리
\ref{moduli-lemma-pic-gerbe-over-pic-functor}와 『스택의 사상』의 보조정리
\ref{stacks-morphisms-lemma-gerbe-fppf}), 『스택의 사상』의 보조정리
\ref{stacks-morphisms-lemma-flat-finite-presentation-permanence}에 의해
$\Picardstack_{X/B} \to B$가 국소적으로 유한 표시임을 보이면 충분하다.
이는 보조정리 \ref{moduli-lemma-pic-qs-lfp}에서 따른다.
\end{proof}

\begin{lemma}
\label{moduli-lemma-pic-functor-uniqueness-part}
이 절의 다른 가정들에 더하여 $X \to B$의 기하적 섬유들이 정역이라고
하자. 그러면 $\Picardfunctor_{X/B} \to B$는 분리이다.
\end{lemma}

\begin{proof}
$\Picardfunctor_{X/B} \to B$가 준분리이므로 값매김 판정법의 유일성
부분을 확인하면 충분하다. 『대수공간의 사상』의 보조정리
\ref{spaces-morphisms-lemma-valuative-criterion-separatedness}를 보라.
이는 곧바로 다음 문제로 귀착된다. 다음이 주어졌다고 하자.
\begin{enumerate}
\item 분수체가 $K$인 값매김환 $R$,
\item $R$ 위에서 고유하고 평탄하며 기하적 섬유가 정역인 대수공간 $X$,
\item $a|_{\Spec(K)} = 0$을 만족하는 원소
$a \in \Picardfunctor_{X/R}(R)$.
\end{enumerate}
이때 $a = 0$임을 증명해야 한다. 전사 평탄 사상
$\Picardstack_{X/R} \to \Picardfunctor_{X/R}$에
『스택의 사상』의 보조정리
\ref{stacks-morphisms-lemma-lift-valuation-ring-through-flat-morphism}을
적용하자(전사성과 평탄성은 보조정리
\ref{moduli-lemma-pic-gerbe-over-pic-functor}와 『스택의 사상』의 보조정리
\ref{stacks-morphisms-lemma-gerbe-fppf}에서 따른다).
$R$을 한 확대환으로 바꾼 뒤에는 $a$가 가역 $\mathcal{O}_X$-가군
$\mathcal{L}$로 주어진다고 가정할 수 있다.
$a|_{\Spec(K)} = 0$이므로 『몫』의 보조정리
\ref{quot-lemma-flat-geometrically-connected-fibres}에 의해
$\mathcal{L}_K \cong \mathcal{O}_{X_K}$이다.

\medskip\noindent
구조 사상을 $f : X \to \Spec(R)$로 표기하자.
$\eta, 0 \in \Spec(R)$를 각각 일반점과 닫힌점이라 하자.
$\Spec(R)$ 위의 완전 복합체
$K = Rf_*\mathcal{L}$과 $M = Rf_*(\mathcal{L}^{\otimes -1})$을 생각하자.
『대수공간의 유도 범주』의 보조정리
\ref{spaces-perfect-lemma-flat-proper-perfect-direct-image-general}를 보라.
$K$와 $M$에 딸린 『대수공간의 유도 범주』의 보조정리
\ref{spaces-perfect-lemma-jump-loci}의 함수
$\beta_{K, i}, \beta_{M, i} : \Spec(R) \to \mathbf{Z}$를 생각하자.
$K$와 $M$의 형성은 밑변환과 가환하므로(위에서 인용한 보조정리를 보라),
『체 위의 대수공간』의 보조정리
\ref{spaces-over-fields-lemma-proper-geometrically-reduced-global-sections}와
$f$의 섬유에 대한 가정에 의해
$\beta_{K, 0}(\eta) = \beta_{M, 0}(\eta) = 1$이다.
위 반연속성에 의해
$\beta_{K, 0}(0) \geq 1$이고 $\beta_{M, 0}(0) \geq 1$이다.
『체 위의 대수공간』의 보조정리
\ref{spaces-over-fields-lemma-characterize-trivial-pic-integral}에 의해
$\mathcal{L}$을 특수 섬유 $X_0$에 제한한 것은 자명하다고 결론 내린다.
그러면 다시 위와 같이
$\beta_{K, 0}(0) = \beta_{M, 0} = 1$이다.
이제 『대수학 더 보기』의 보조정리
\ref{more-algebra-lemma-lift-pseudo-coherent-from-residue-field}에 의해
$K$를 다음 형태의 복합체로 나타낼 수 있다.
$$
\ldots \to 0 \to R \to R^{\oplus \beta_{K, 1}(0)} \to
R^{\oplus \beta_{K, 2}(0)} \to \ldots
$$
$\beta_{K, 0}(\eta) = 1$이므로
$R \to R^{\oplus \beta_{K, 1}(0)}$은 영이다. 다시 말해
$D(R)$에서 $K = R \oplus \tau_{\geq 1}(K)$이고, 어떤
$b \in \mathbf{Z}$에 대해 $\tau_{\geq 1}(K)$의 토르 진폭은
$[1, b]$ 안에 있다. 따라서 $X_0$로의 제한 $s_0$이 소멸하지 않는
대역 단면 $s \in H^0(X, \mathcal{L})$이 존재한다
($K$의 형성이 밑변환과 가환하기 때문이다).
그러면 $s : \mathcal{O}_X \to \mathcal{L}$은 그 $X_0$로의 제한이
동형인 가역층들의 사상이므로, 원하는 대로 그 자체가 동형이다.
\end{proof}

\begin{lemma}
\label{moduli-lemma-pic-functor-curves-smooth}
이 절의 다른 가정들에 더하여 $f : X \to B$의 상대 차원이
$\leq 1$이라고 하자. 그러면 $\Picardfunctor_{X/B} \to B$는 매끄럽다.
\end{lemma}

\begin{proof}
보조정리 \ref{moduli-lemma-pic-curves-smooth}에 의해
$\Picardstack_{X/B} \to B$가 매끄러움을 안다. 보조정리
\ref{moduli-lemma-pic-gerbe-over-pic-functor}와 『스택의 사상』의 보조정리
\ref{stacks-morphisms-lemma-gerbe-smooth}를 결합하면 사상
$\Picardstack_{X/B} \to \Picardfunctor_{X/B}$는 전사이고 매끄럽다.
따라서 $U$가 스킴이고 $U \to \Picardstack_{X/B}$가 전사이고
매끄러우면, $U \to \Picardfunctor_{X/B}$는 전사이고 매끄러우며
$U \to B$도 전사이고 매끄럽다(이 성질들은 합성에 의해 보존되기
때문이다). 그러므로 예를 들어 『대수공간 위의 내림』의 보조정리
\ref{spaces-descent-lemma-syntomic-smooth-etale-permanence}에 의해
$\Picardfunctor_{X/B} \to B$는 매끄럽다.
\end{proof}






\section{상대 사상의 성질}
\label{moduli-section-relative-morphisms}

\noindent
$B$를 대수공간이라 하자. $X$와 $Y$를 $B$ 위의 대수공간이라 하고,
$Y \to B$는 평탄하고 고유이며 유한 표시이고 $X \to B$는 분리이고
유한 표시라고 하자. 그러면 상대 사상의 함자
$\mathit{Mor}_B(Y, X)$는 $B$ 위에서 국소적으로 유한 표시인 대수공간이다.
『몫』의 명제 \ref{quot-proposition-Mor}를 보라.

\begin{lemma}
\label{moduli-lemma-Mor-diagonal-closed}
$\mathit{Mor}_B(Y, X) \to B$의 대각사상은 유한 표시인 닫힌 몰입이다.
\end{lemma}

\begin{proof}
열린 몰입
$\mathit{Mor}_B(Y, X) \to \Hilbfunctor_{Y \times_B X/B}$가 존재한다.
『몫』의 보조정리 \ref{quot-lemma-Mor-into-Hilb-open}을 보라.
따라서 이 보조정리는 보조정리
\ref{moduli-lemma-hilb-diagonal-closed}에서 따른다.
\end{proof}

\begin{lemma}
\label{moduli-lemma-Mor-s-lfp}
사상 $\mathit{Mor}_B(Y, X) \to B$는 분리이고 국소적으로 유한 표시이다.
\end{lemma}

\begin{proof}
$\mathit{Mor}_B(Y, X) \to B$가 분리임을 확인하려면 그 대각사상이
닫힌 몰입임을 보여야 한다. 이는 보조정리
\ref{moduli-lemma-Mor-diagonal-closed}에 의해 참이다. 두 번째 명제는
『몫』의 명제 \ref{quot-proposition-Mor}의 일부이다.
\end{proof}

\begin{lemma}
\label{moduli-lemma-Isom-in-Mor}
$B, X, Y$를 이 절의 도입부에서와 같이 두고, 추가로 $X \to B$가
고유라고 하자. 그러면 동형들로 이루어진 부분함자
$\mathit{Isom}_B(Y, X) \subset \mathit{Mor}_B(Y, X)$는 열린 부분공간이다.
\end{lemma}

\begin{proof}
『대수공간의 사상 더 보기』의 보조정리
\ref{spaces-more-morphisms-lemma-where-isomorphism}에서 즉시 따른다.
\end{proof}

\begin{remark}[수치 불변량]
\label{moduli-remark-Mor-numerical}
$B, X, Y$를 이 절의 도입부에서와 같이 두자. $I$를 집합이라 하고
각 $i \in I$에 대해 $E_i \in D(\mathcal{O}_{Y \times_B X})$가
완전하다고 하자. $P : I \to \mathbf{Z}$를 함수라 하자.
$$
\mathit{Mor}_B(Y, X) \subset
\Hilbfunctor_{Y \times_B X/B}
$$
가 열린 부분공간임을 상기하자. 『몫』의 보조정리
\ref{quot-lemma-Mor-into-Hilb-open}을 보라. 따라서
$$
\mathit{Mor}^P_B(Y, X) =
\mathit{Mor}_B(Y, X) \cap \Hilbfunctor^P_{Y \times_B X/B}
$$
로 정의할 수 있으며, 여기서 $\Hilbfunctor^P_{Y \times_B X/B}$는
주석 \ref{moduli-remark-hilb-numerical}에서와 같다. 사상
$$
\mathit{Mor}^P_B(Y, X) \longrightarrow \mathit{Mor}_B(Y, X)
$$
은 평탄한 닫힌 몰입이며, 예를 들어 $I$가 유한이거나 $B$가 국소
뇌터이거나 어떤 가역 $\mathcal{O}_{Y \times_B X}$-가군
$\mathcal{L}$에 대해 $I = \mathbf{Z}$이고
$E_i = \mathcal{L}^{\otimes i}$이면 열린 동시에 닫힌 몰입이다.
마지막 경우에는 때때로 $\mathit{Mor}^{P, \mathcal{L}}_B(Y, X)$라는
표기를 쓴다.
\end{remark}

\begin{lemma}
\label{moduli-lemma-Mor-qc-over-base}
$B, X, Y$를 이 절의 도입부에서와 같이 두고, $\mathcal{L}$은
$X/B$ 위에서 풍부하고 $\mathcal{N}$은 $Y/B$ 위에서 풍부하다고 하자.
『대수공간 위의 인자』의 정의
\ref{spaces-divisors-definition-relatively-ample}를 보라.
$P$를 수치 다항식이라 하자. 그러면
$$
\mathit{Mor}^{P, \mathcal{M}}_B(Y, X) \longrightarrow B
$$
는 분리이고 유한 표시이며, 여기서
$\mathcal{M} = \text{pr}_1^*\mathcal{N}
\otimes_{\mathcal{O}_{Y \times_B X}} \text{pr}_2^*\mathcal{L}$이다.
\end{lemma}

\begin{proof}
보조정리 \ref{moduli-lemma-Mor-s-lfp}에 의해 사상
$\mathit{Mor}_B(Y, X) \to B$는 분리이고 국소적으로 유한 표시이다.
따라서 주석 \ref{moduli-remark-Mor-numerical}의 열린 동시에 닫힌 부분공간
$\mathit{Mor}^{P, \mathcal{M}}_B(Y, X)$가 $B$ 위에서 준콤팩트임을
보이면 충분하다.

\medskip\noindent
이 문제는 $B$ 위에서 에탈 국소적이다
(『대수공간의 사상』의 보조정리
\ref{spaces-morphisms-lemma-quasi-compact-local}). 따라서 $B$가
아핀이라고 가정할 수 있다.

\medskip\noindent
$B = \Spec(\Lambda)$라 하자. $X$와 $Y$는 스킴이고
$\mathcal{L}$과 $\mathcal{N}$은 각각 $X$와 $Y$ 위의 풍부한 가역층임에
유의하자(이는 정의에서 즉시 따른다).
$\Lambda = \colim \Lambda_i$를 그 유한형 $\mathbf{Z}$-부분대수들의
쌍대극한으로 쓰자. 그러면 어떤 $i$에 대해
$B_i = \Spec(\Lambda_i)$ 위에서 보조정리의 조건을 만족하는 계
$X_i, Y_i, \mathcal{L}_i, \mathcal{N}_i$를 찾아, 이를 $B$로 밑변환하면
$X, Y, \mathcal{L}, \mathcal{N}$이 되게 할 수 있다. 이는 『극한』의
보조정리
\ref{limits-lemma-descend-finite-presentation}($X_i$, $Y_i$를 찾는 데),
\ref{limits-lemma-descend-invertible-modules}
($\mathcal{L}_i$, $\mathcal{N}_i$를 찾는 데),
\ref{limits-lemma-descend-separated-finite-presentation}
($X_i \to B_i$를 분리로 만드는 데),
\ref{limits-lemma-eventually-proper}($Y_i \to B_i$를 고유로 만드는 데),
그리고 \ref{limits-lemma-limit-ample}
($\mathcal{L}_i$, $\mathcal{N}_i$를 풍부하게 만드는 데)에서 따른다.
다음 등식과 $\mathit{Mor}^P_B(Y, X)$에 대한 같은 등식에 의해
다음 문단에서 다루는 경우로 귀착된다.
$$
\mathit{Mor}_B(Y, X) = B \times_{B_i} \mathit{Mor}_{B_i}(Y_i, X_i)
$$

\medskip\noindent
$B$가 뇌터 아핀 스킴이라고 하자. 『성질』의 보조정리
\ref{properties-lemma-ample-on-product}에 의해 $\mathcal{M}$은 풍부하다.
보조정리 \ref{moduli-lemma-hilb-qc-over-base}에 의해
$\Hilbfunctor^{P, \mathcal{M}}_{Y \times_B X/B}$는 $B$ 위에서 유한
표시이고, 따라서 뇌터이다. 구성에 의해
$$
\mathit{Mor}^{P, \mathcal{M}}_B(Y, X) =
\mathit{Mor}_B(Y, X) \cap
\Hilbfunctor^{P, \mathcal{M}}_{Y \times_B X/B}
$$
는 $\Hilbfunctor^{P, \mathcal{M}}_{Y \times_B X/B}$의 열린 부분공간이고,
따라서 준콤팩트이다(뇌터 대수공간의 열린 부분공간은 준콤팩트이므로).
\end{proof}








\section{편극된 고유 스킴의 스택의 성질}
\label{moduli-section-polarized}

\noindent
이 절에서는 모듈라이 스택
$$
\Polarizedstack \longrightarrow \Spec(\mathbf{Z})
$$
의 성질을 논한다. 스킴 $S$ 위의 이 스택의 단면 범주는 상대적으로
풍부한 가역층이 주어진, $S$ 위에서 고유하고 평탄하며 유한 표시인
스킴의 범주이다. 이는 『몫』의 정리
\ref{quot-theorem-polarized-algebraic}에 의해 대수적 스택이다.

\begin{lemma}
\label{moduli-lemma-polarized-diagonal-separated-fp}
$\Polarizedstack$의 대각사상은 분리이고 유한 표시이다.
\end{lemma}

\begin{proof}
$\Polarizedstack$가 극한을 보존하는 대수적 스택임을 상기하자.
『몫』의 보조정리 \ref{quot-lemma-polarized-limits}를 보라.
『스택의 극한』의 보조정리
\ref{stacks-limits-lemma-limit-preserving-diagonal}에 의해 이는
$\Delta : \Polarizedstack \to \Polarizedstack \times \Polarizedstack$가
극한을 보존함을 뜻한다. 따라서 『스택의 극한』의 명제
\ref{stacks-limits-proposition-characterize-locally-finite-presentation}에
의해 $\Delta$는 국소적으로 유한 표시이다.

\medskip\noindent
$\Delta$가 분리임을 증명하자. 이를 위해서는 아핀 스킴 $U$와
$U$ 위의 $\Polarizedstack$의 두 대상
$\upsilon = (Y, \mathcal{N})$과 $\chi = (X, \mathcal{L})$이 주어졌을 때
대수공간
$$
\mathit{Isom}_{\Polarizedstack}(\upsilon, \chi)
$$
가 분리임을 보이면 충분하다. 동형 $\upsilon_T \to \chi_T$에 그
기저 동형 $Y_T \to X_T$를 대응시키는 규칙은 사상
$$
\mathit{Isom}_{\Polarizedstack}(\upsilon, \chi)
\longrightarrow
\mathit{Isom}_U(Y, X)
$$
을 정의한다. 보조정리 \ref{moduli-lemma-Mor-s-lfp}와
\ref{moduli-lemma-Isom-in-Mor}에서 표적이 분리 대수공간임을 보았으므로,
이 사상이 분리임을 증명하면 충분하다. 어떤 스킴 $T/U$ 위의 동형
$f : Y_T \to X_T$가 주어지면 명백히
$$
\mathit{Isom}_{\Polarizedstack}(\upsilon, \chi)
\times_{\mathit{Isom}_U(Y, X), [f]} T
=
\mathit{Isom}(\mathcal{N}_T, f^*\mathcal{L}_T)
$$
이다. 여기서 $[f] : T \to \mathit{Isom}_U(Y, X)$는 $f$에 대응하는
$T$-값 점을 나타내고,
$\mathit{Isom}(\mathcal{N}_T, f^*\mathcal{L}_T)$는 절
\ref{moduli-section-hom-isom}에서 논한 대수공간이다. 이 대수공간은 $U$ 위에서
아핀이므로, 이 주장은 $\Delta$가 분리임을 함의한다.

\medskip\noindent
증명을 마치기 위해 $\Delta$가 준콤팩트임을 보이자. $\Delta$는
대수공간으로 표현가능하므로, 전사 매끄러운 사상
$U \to \Polarizedstack \times \Polarizedstack$에 의한 $\Delta$의
밑변환이 준콤팩트임을 확인하면 충분하다(예를 들어 『스택의 성질』의
보조정리 \ref{stacks-properties-lemma-check-property-covering}를 보라).
$U = \coprod U_i$가 아핀 열린 부분들의 서로소 합이라고 가정할 수 있다.
$\Polarizedstack$는 극한을 보존하므로(위를 보라),
$\Polarizedstack \to \Spec(\mathbf{Z})$는 국소적으로 유한 표시이고,
따라서 $U_i \to \Spec(\mathbf{Z})$도 국소적으로 유한 표시이다
(『스택의 극한』의 명제
\ref{stacks-limits-proposition-characterize-locally-finite-presentation}과
『스택의 사상』의 보조정리
\ref{stacks-morphisms-lemma-composition-finite-presentation} 및
\ref{stacks-morphisms-lemma-smooth-locally-finite-presentation}).
특히 $U_i$는 뇌터 아핀이다. 이로써 다음 문단에서 다루는 경우로
귀착된다.

\medskip\noindent
이 문단에서는 뇌터 아핀 스킴 $U$와 $U$ 위의 $\Polarizedstack$의
두 대상 $\upsilon = (Y, \mathcal{N})$과
$\chi = (X, \mathcal{L})$이 주어졌을 때 대수공간
$$
\mathit{Isom}_{\Polarizedstack}(\upsilon, \chi)
$$
가 준콤팩트임을 보인다. $U$의 연결 성분들은 열린 동시에 닫혀 있으므로
$U$를 이들로 바꿀 수 있다. 따라서 $U$가 연결이라고 가정하며 실제로
그렇게 하겠다. $u \in U$를 한 점이라 하자.
$n \mapsto \chi(Y_u, \mathcal{N}_u^{\otimes n})$을 힐베르트 다항식
$P$라 하자. 『다양체』의 보조정리
\ref{varieties-lemma-numerical-polynomial-from-euler}를 보라.
$U$가 연결이고 함수
$u \mapsto \chi(Y_u, \mathcal{N}_u^{\otimes n})$들이 국소 상수이므로
(『스킴의 유도 범주』의 보조정리
\ref{perfect-lemma-chi-locally-constant-geometric}를 보라),
$U$의 모든 점에서 같은 힐베르트 다항식을 얻는다.
$Y \times_U X$ 위에서
$\mathcal{M} = \text{pr}_1^*\mathcal{N}
\otimes_{\mathcal{O}_{Y \times_U X}} \text{pr}_2^*\mathcal{L}$로 두자.
$U$ 위의 어떤 스킴 $T$에 대해
$(f, \varphi) \in \mathit{Isom}_{\Polarizedstack}(\upsilon, \chi)(T)$가
주어지면 모든 $t \in T$에 대해
$$
\chi(Y_t, (\text{id} \times f)^*\mathcal{M}^{\otimes n}) =
\chi(Y_t,
\mathcal{N}_t^{\otimes n} \otimes_{\mathcal{O}_{Y_t}}
f_t^*\mathcal{L}_t^{\otimes n}) =
\chi(Y_t, \mathcal{N}_t^{\otimes 2n}) = P(2n)
$$
이다. 가운데 등식에서는 동형
$\varphi : f^*\mathcal{L}_T \to \mathcal{N}_T$를 사용한다.
$P'(t) = P(2t)$로 두면 앞서 본 사상
$$
\mathit{Isom}_{\Polarizedstack}(\upsilon, \chi)
\longrightarrow
\mathit{Isom}_U(Y, X)
$$
의 상은 교집합
$$
\mathit{Isom}_U(Y, X) \cap \mathit{Mor}^{P', \mathcal{M}}_U(Y, X)
$$
에 포함된다. 이 교집합은 $\mathit{Mor}_U(Y, X)$의 열린 부분공간들의
교집합이다(보조정리 \ref{moduli-lemma-Isom-in-Mor}와 주석
\ref{moduli-remark-Mor-numerical}을 보라).
이제 보조정리 \ref{moduli-lemma-Mor-qc-over-base}에 의해
$\mathit{Mor}^{P', \mathcal{M}}_U(Y, X)$는 $U$ 위에서 유한 표시인
뇌터 대수공간이다. 따라서 이 교집합도 뇌터 대수공간이다. 사상
$$
\mathit{Isom}_{\Polarizedstack}(\upsilon, \chi)
\longrightarrow
\mathit{Isom}_U(Y, X) \cap \mathit{Mor}^{P', \mathcal{M}}_U(Y, X)
$$
이 아핀이므로(위를 보라) 결론을 얻는다.
\end{proof}

\begin{lemma}
\label{moduli-lemma-polarized-qs-lfp}
사상 $\Polarizedstack \to \Spec(\mathbf{Z})$는 준분리이고 국소적으로
유한 표시이다.
\end{lemma}

\begin{proof}
$\Polarizedstack \to \Spec(\mathbf{Z})$가 준분리임을 확인하려면 그
대각사상이 준콤팩트이고 준분리임을 보여야 한다. 이는 보조정리
\ref{moduli-lemma-polarized-diagonal-separated-fp}에서 즉시 따른다.
$\Polarizedstack \to \Spec(\mathbf{Z})$가 국소적으로 유한 표시임을
증명하려면 $\Polarizedstack$가 극한을 보존함을 보이면 충분하다.
『스택의 극한』의 명제
\ref{stacks-limits-proposition-characterize-locally-finite-presentation}를
보라. 이는 『몫』의 보조정리 \ref{quot-lemma-polarized-limits}이다.
\end{proof}

\begin{lemma}
\label{moduli-lemma-bounded-polarized}
$n \geq 1$을 정수, $P$를 수치 다항식이라 하자. 부분집합
$$
T \subset |\Polarizedstack|
$$
이 다음 성질을 만족한다고 하자. 모든 $\xi \in T$에 대해 체 $k$와
$\xi$를 나타내는 $k$ 위의 $\Polarizedstack$의 대상
$(X, \mathcal{L})$이 존재하여 다음을 만족한다.
\begin{enumerate}
\item $X$ 위의 $\mathcal{L}$의 힐베르트 다항식은 $P$이고,
\item $i^*\mathcal{O}_{\mathbf{P}^n}(1) \cong \mathcal{L}$을 만족하는
닫힌 몰입 $i : X \to \mathbf{P}^n_k$가 존재한다.
\end{enumerate}
그러면 $T$는 뇌터 위상공간이고, 특히 준콤팩트이다.
\end{lemma}

\begin{proof}
$|\Polarizedstack|$는 국소 뇌터 위상공간임에 유의하자.
『스택의 사상』의 보조정리
\ref{stacks-morphisms-lemma-Noetherian-topology}를 보라
(여기서는 $\Spec(\mathbf{Z})$가 뇌터라는 사실, 따라서 보조정리
\ref{moduli-lemma-polarized-qs-lfp}와 『스택의 사상』의 보조정리
\ref{stacks-morphisms-lemma-locally-finite-type-locally-noetherian}에 의해
$\Polarizedstack$가 국소 뇌터 대수적 스택이라는 사실도 사용한다).
따라서 $|\Polarizedstack|$의 임의의 준콤팩트 부분집합은 뇌터
위상공간이고, 그러한 집합의 임의의 부분집합 역시 뇌터이다.
『위상수학』의 보조정리 \ref{topology-lemma-finite-union-Noetherian}과
\ref{topology-lemma-Noetherian}을 보라. 따라서 $T$를 포함하는 준콤팩트
부분집합을 하나 찾기만 하면 된다.

\medskip\noindent
보조정리 \ref{moduli-lemma-hilb-proper-over-base}에 의해 대수공간
$$
H =
\Hilbfunctor^{P, \mathcal{O}(1)}_{\mathbf{P}^n_\mathbf{Z}/\Spec(\mathbf{Z})}
$$
는 $\Spec(\mathbf{Z})$ 위에서 고유이다. 『몫』의 보조정리
\ref{quot-lemma-extend-hilb-to-spaces}\footnote{나중에
(여기에 미래의 참조를 삽입하라) $H$가 스킴임을 볼 것이므로, 이
보조정리와 『몫』의 보조정리
\ref{quot-lemma-extend-polarized-to-spaces}를 사용할 필요가 없다.}에
의해 $H$의 항등 사상은 닫힌 부분공간
$$
Z \subset \mathbf{P}^n_H
$$
에 대응한다. 이 부분공간은 $H$ 위에서 고유하고 평탄하며 유한 표시이고,
제한 $\mathcal{N} = \mathcal{O}(1)|_Z$는 $Z/H$ 위에서 상대적으로
풍부하며 $Z \to H$의 섬유들 위에서 힐베르트 다항식 $P$를 갖는다.
특히 쌍 $(Z \to H, \mathcal{N})$은 사상
$$
H \longrightarrow \Polarizedstack
$$
을 정의한다. 이 사상은 스킴의 사상 $U \to H$를 족
$(Z_U \to U, \mathcal{N}_U)$의 분류 사상으로 보낸다. 『몫』의
보조정리 \ref{quot-lemma-extend-polarized-to-spaces}를 보라.
$H$는 뇌터 대수공간이므로($\mathbf{Z})$ 위에서 고유이므로),
$|H|$는 뇌터이고 따라서 준콤팩트이다. 사상
$$
|H| \longrightarrow |\Polarizedstack|
$$
은 연속이므로 그 상은 준콤팩트이다. 따라서 $T$가
$|H| \to |\Polarizedstack|$의 상에 포함됨을 증명하면 충분하다.
그러나 가정 (1)과 (2)는 정확히 이것이 성립한다는 사실을 표현한다.
$i^*\mathcal{O}_{\mathbf{P}^n}(1) \cong \mathcal{L}$을 만족하는 닫힌
몰입 $i : X \to \mathbf{P}^n_k$를 어떻게 택하더라도 $H$의 모듈라이
해석에 의해 $H$의 $k$-값 점을 얻는다. 이로써 보조정리의 증명이 끝난다.
\end{proof}







\section{고유 사상 위의 복합체의 모듈라이의 성질}
\label{moduli-section-complexes}

\noindent
$f : X \to B$를 고유하고 평탄하며 유한 표시인 대수공간의 사상이라
하자. 그러면 음의 자기 Ext가 소멸하는 상대 완전 복합체를 매개화하는
스택 $\Complexesstack_{X/B}$는 대수적이다. 『몫』의 정리
\ref{quot-theorem-complexes-algebraic}를 보라.

\begin{lemma}
\label{moduli-lemma-complexes-diagonal-affine-fp}
$B$ 위에서 $\Complexesstack_{X/B}$의 대각사상은 아핀이며 유한 표시이다.
\end{lemma}

\begin{proof}
대각사상이 대수공간으로 표현가능함은 『몫』의 보조정리
\ref{quot-lemma-complexes-diagonal}에서 보였다. 그 증명을 살펴보면
다음을 보여야 함을 알 수 있다. $B$ 위의 스킴 $T$와 대상
$E, E' \in D(\mathcal{O}_{X_T})$가 주어져서
$(T, E)$와 $(T, E')$가 $T$ 위의 $\Complexesstack_{X/B}$의 섬유 범주의
대상이라고 하자. 그러면
$\mathit{Isom}(E, E') \to T$는 아핀이며 유한 표시이다.
여기서 $\mathit{Isom}(E, E')$는 다음 함자이다.
$$
(\Sch/T)^{opp} \to \textit{Sets},\quad
T' \mapsto \{\varphi : E_{T'} \to E'_{T'}
\text{ isomorphism in }D(\mathcal{O}_{X_{T'}})\}
$$
여기서 $E_{T'}$와 $E'_{T'}$는 각각 $E$와 $E'$를 $X_{T'}$로 유도
당김한 것이다. 다음 규칙으로 정의되는 함자
$H = \SheafHom(E, E')$를 생각하자.
$$
(\Sch/T)^{opp} \to \textit{Sets},\quad
T' \mapsto \Hom_{\mathcal{O}_{X_{T'}}}(E_T, E'_T)
$$
『몫』의 보조정리
\ref{quot-lemma-complexes-open-neg-exts-vanishing}에 의해 이는 $T$ 위에서
아핀이고 유한 표시인 대수공간이다.
$H' = \SheafHom(E', E)$, $I = \SheafHom(E, E)$, 그리고
$I' = \SheafHom(E', E')$에 대해서도 마찬가지이다. 따라서
$$
\mathit{Isom}(E, E') = (H' \times_T H) \times_{c, I \times_T I', \sigma} T
$$
임을 알 수 있다. 여기서
$c(\varphi', \varphi) = (\varphi \circ \varphi', \varphi' \circ \varphi)$이고
$\sigma = (\text{id}, \text{id})$이다(『몫』의 명제
\ref{quot-proposition-isom}의 증명과 비교하라). 따라서
$\mathit{Isom}(E, E')$는 $T$ 위의 아핀 스킴들의 섬유곱으로서
$T$ 위에서 아핀이다. 마찬가지로 $\mathit{Isom}(E, E')$는
$T$ 위에서 유한 표시이다.
\end{proof}

\begin{lemma}
\label{moduli-lemma-complexes-qs-lfp}
사상 $\Complexesstack_{X/B} \to B$는 준분리이고 국소적으로
유한 표시이다.
\end{lemma}

\begin{proof}
$\Complexesstack_{X/B} \to B$가 준분리임을 확인하려면 그 대각사상이
준콤팩트이고 준분리임을 보여야 한다. 이는 보조정리
\ref{moduli-lemma-complexes-diagonal-affine-fp}에서 즉시 따른다.
$\Complexesstack_{X/B} \to B$가 국소적으로 유한 표시임을 증명하려면
$\Complexesstack_{X/B} \to B$가 극한을 보존함을 보여야 한다.
『스택의 극한』의 명제
\ref{stacks-limits-proposition-characterize-locally-finite-presentation}를
보라. 이는 『몫』의 보조정리 \ref{quot-lemma-complexes-limits}에서 따른다
(작은 세부사항은 생략한다).
\end{proof}
